\documentclass[10pt,a4paper]{article}
\usepackage[margin=18mm,headheight=14pt]{geometry}
\usepackage{fontspec}
\setmainfont{Latin Modern Roman}
\usepackage{amsmath,amssymb,mathtools,enumitem,graphicx,longtable,array,fancyhdr}
\usepackage{tikz-cd}
\usepackage[unicode,hidelinks]{hyperref}
\providecommand{\GL}{\operatorname{GL}}
\providecommand{\SL}{\operatorname{SL}}
\providecommand{\PGL}{\operatorname{PGL}}
\providecommand{\Hom}{\operatorname{Hom}}
\providecommand{\Iso}{\operatorname{Iso}}
\providecommand{\Isom}{\operatorname{Isom}}
\providecommand{\Ker}{\operatorname{Ker}}
\providecommand{\Aut}{\operatorname{Aut}}
\providecommand{\Lie}{\operatorname{Lie}}
\providecommand{\Sym}{\operatorname{Sym}}
\providecommand{\Tr}{\operatorname{Tr}}
\providecommand{\Ind}{\operatorname{Ind}}
\providecommand{\car}{\operatorname{car}}
\providecommand{\im}{\operatorname{im}}
\providecommand{\red}{\operatorname{red}}
\providecommand{\sprep}{\operatorname{sp}}
\providecommand{\sprod}{\mathop{\prod}\nolimits'}
\providecommand{\R}{\mathbb R}
\providecommand{\Q}{\mathbb Q}
\providecommand{\Z}{\mathbb Z}
\providecommand{\C}{\mathbb C}
\providecommand{\A}{\mathbb A}
\providecommand{\Af}{\mathbb A^f}
\providecommand{\GA}{G_{\mathbb A}}
\providecommand{\SG}{SG}
\providecommand{\OO}{\mathcal O}
\providecommand{\M}{\mathrm M}
\providecommand{\calR}{\mathcal R}
\providecommand{\calJ}{\mathfrak F}
\providecommand{\calH}{\mathcal H}
\providecommand{\calK}{\mathcal K}
\providecommand{\cK}{\mathcal K}
\providecommand{\cS}{\mathcal S}
\providecommand{\iso}{\xrightarrow{\sim}}
\setlength{\parindent}{1em}
\setlength{\parskip}{3pt}
\setlist{nosep,leftmargin=2em}
\allowdisplaybreaks
\emergencystretch=2em
\pagestyle{fancy}
\fancyhf{}
\fancyhead[L]{\small P. Deligne --- D017}
\fancyhead[R]{\small\nouppercase{\leftmark}}
\fancyfoot[C]{\thepage}
\begin{document}
\begin{center}\Large Modular forms and representations of GL(2)\\[5pt]\large P. Deligne\\[5pt] Standalone English edition\end{center}
\noindent\small Authority: 51 physical pages, printed pages 55--105. Source spelling, abbreviations and apparent errors are retained. Page labels below are editorial anchors, not source text. Typography is reflowed; source folios and running identifiers are excluded. The separate apparatus records authority restorations.\normalsize\par

\clearpage
\markboth{Authority 1 / printed 55}{}
\hypertarget{authority-1}{}
\begin{center}
\underline{MODULAR FORMS AND REPRESENTATIONS OF GL(2)}

\bigskip
\underline{by P. Deligne}
\end{center}
\vspace{2em}
\noindent\begin{tabular}{@{}p{.87\linewidth}r@{}}
Introduction&2\\[.6em]
\underline{§ 0. Preliminaries}&\\
\quad 0.0 Notation&4\\
\quad 0.1 Adelic lattices&5\\
\quad 0.2 Admissible representations&7\\[.6em]
\underline{§ 1. Functions of lattices}&\\
\quad 1.1 Functions on $\GL(2,\R)$&8\\
\quad 1.2 Functions on $\GL(2,\A)$&15\\
\quad 1.3 Cusps&20\\[.6em]
\underline{§ 2. Modular forms and the spectrum of $\GL_2$}&\\
\quad 2.1 Holomorphic modular forms and\newline\hspace*{2em}representations of $\GL(2,\R)$&22\\
\quad 2.2 Kirillov model and new vector&24\\
\quad 2.3 Kirillov model and Hecke operators&29\\
\quad 2.4 Atkin-Lehner newforms&30\\
\quad 2.5 Expansions at the cusps&34\\[.6em]
\underline{§ 3. Local class field theory and representations of $\GL(2,F)$}&\\
\quad 3.1 Preliminaries&38\\
\quad 3.2 Representations of $\GL(2,K)$&41\\[.6em]
Bibliography&51
\end{tabular}

\clearpage
\markboth{Authority 2 / printed 56}{}
\hypertarget{authority-2}{}
\subsection*{Introduction}
This exposition, which makes no claim to originality, is intended as a complement to Robert's report [13] on Jacquet-Langlands [9]. It can be divided into two parts rather different in spirit. Sections 0, 1, and 2 are intended as a bridge, in the theory of modular forms, between the classical point of view and that of representations of $\GL(2)$. Section 3 is a compendium of results from the theory of representations of $\GL(2,K)$ ($K$ a local field). It depends on Sections 2, 3, and 8 of [4] (in this volume), which themselves are independent of the rest of that article.

The fundamental global theorems will not be discussed here (relations among the spectrum of $\GL(2)$, the spectrum of quaternion algebras, and Dirichlet series with functional equations); for these we refer to Robert's cited report and, of course, to [9]. We shall dwell on the following points.

(A) The relation between the classical point of view: modular forms as functions on the Poincaré half-plane, or as functions of lattices, and the representation-theoretic point of view, in which one seeks to decompose $L_0^2(\GL(2,\A)/\GL(2,\Q))$ as a sum of irreducible representations of the adelic group.

(B) In Section 2, after the minimum required concerning $\GL(2,\R)$, we give two corollaries of the existence-and-uniqueness theorem for the Kirillov model of irreducible admissible representations of $\GL(2,K)$ ($K$ a nonarchimedean local field).

a) The theory of the new forms of Atkin and Lehner [1]: the proof given here is simpler, and closer to Atkin-Lehner's, than Casselman's [3], but goes less far.

b) The multiplicity-one theorem: if, for almost every prime $p$, $\pi_p$ is an irreducible admissible representation of $\GL(2,\Q_p)$, there exists at most one irreducible admissible subrepresentation of $\GL(2,\A)$ in $L_0^2(\GL(2,\A)/\GL(2,\Q))$ whose local components agree with the $\pi_p$ at almost all primes and whose local component at infinity is prescribed. The proof is very elementary and does not use global $L$-functions.

\clearpage
\markboth{Authority 3 / printed 57}{}
\hypertarget{authority-3}{}
(C) Section 3 is a compendium of results on the relations between irreducible admissible representations of $\GL(2,K)$ ($K$ a local field) and two-dimensional representations of the Galois group (or rather the Weil group) of $K$. For each theorem cited, I have tried, insofar as possible, to give a reference in which the proof requires a minimum of calculation. This section does not discuss either the explicit construction of the representations (for which one may consult Cartier's exposition [2] in this volume) or the decomposition of their restrictions to the maximal compact subgroup (see Silberger [14]).

As is apparent from the preceding, I have considered only $\GL(2,\Q)$, not $\GL(2,F)$ for a global field $F$. In view of the emphasis placed on local theory, this has little consequence. I have also restricted myself to holomorphic modular forms. The general case is treated in the same way once the results on $\GL(2,\R)$ (and $\GL(2,\C)$) proved in [9], §§5 and 6, or [15], Chapter VIII, are available.

Among the many questions passed over in silence, I also mention

\noindent - Eisenstein series, for which one may refer to Hecke [8], 24;

\noindent - the relation between $\SL(2)$ and $\GL(2)$, for which I refer to [10].

Besides the pervasive influence of [9], I may cite, as sources of which I am particularly conscious:

for § 1, no. 2, the introduction to [11]; for § 3, [12]; for § 2, no. 4, [3].

\clearpage
\markboth{Authority 4 / printed 58}{}
\hypertarget{authority-4}{}
\section*{0. Preliminaries}
\subsection*{0.0. Notation}
Set
\[
\begin{aligned}
\widehat{\Z}&=\varprojlim\Z/n\Z &&(n\in\mathbb N^+,\text{ ordered by divisibility})\\
\Z_p&=\varprojlim_n\Z/p^n\Z &&\text{(the same limit over a subset of the indices)}\\
\A^f&=\widehat{\Z}\otimes\Q=\bigcup\frac1n\widehat{\Z}\\
\Q_p&=\Z_p\otimes\Q=\bigcup\frac1{p^n}\widehat{\Z}_p\\
\A&=\R\times\A^f=(\R\times\widehat{\Z})\otimes\Q
\end{aligned}
\]
We have
\[
\widehat{\Z}\xrightarrow{\sim}\prod_p\Z_p
\]
and the adele ring $\A$ is the restricted product, relative to the compact open subgroups $\Z_p\subset\Q_p$, of the completions $\Q_v$ ($\R=\Q_\infty$, and the $\Q_p$) of $\Q$.

\noindent $v$ or $v_p:\Q_p\longrightarrow\Z\cup\{\infty\}$: for $x\in p^n\Z_p^*$, $v_p(x)=n$; $v_p(0)=\infty$

\noindent $\|x\|:\Q_p\longrightarrow\R:\|x\|=p^{-v_p(x)}$

\noindent $\|x\|:\A\longrightarrow\R:\|x\|=\prod_v\|x_v\|$, where $x_v\in\Q_v$ are the components of $x$, almost all lying in $\Z_p\subset\Q_p$.

For any ring $A$ and any free $A$-module $V$, set

\noindent $\GL_A(V)$ = the group of $A$-linear automorphisms of $V$, and

\noindent $\GL(n,A)=\GL_A(A^n)$: invertible $n\times n$ matrices.

\noindent local field: a (commutative) locally compact nondiscrete field, for example $\R$, $\C$, or $\Q_p$. We denote such fields by $F$ or $K$, according to the section.

\noindent archimedean local field: $\R$ or $\C$

\noindent nonarchimedean local field: any other local field. For such a field $F$, let $\mathcal O$ denote its valuation ring (a complete discrete valuation ring with finite residue field $\mathbb F_q$), $\pi$ a uniformizer, and $v$ the

\clearpage
\markboth{Authority 5 / printed 59}{}
\hypertarget{authority-5}{}
valuation ($v(\pi)=1$); write $\|\ \|$ for the absolute value $\|x\|=q^{-v(x)}$.

\subsection*{0.1. Adelic lattices}
(0.1.1) Recall that, just as $\Z$ is discrete with compact quotient in $\R$, $\Q$ is discrete with compact quotient in $\A$.

Let $M$ be a free $\A$-module of rank $n$. A \underline{lattice} in $M$ is a $\Q$-vector subspace $R\subset M$ such that
\[
R\otimes_{\Q}\A\xrightarrow{\sim}M.
\]
Equivalently, it is a $\Q$-vector subspace $R\subset M$ that is discrete and has compact quotient (just as lattices in $\R^n$ are discrete subgroups with compact quotient).

Let $R_0\subset\R^n$ be a lattice and, for each prime $\ell$, let $\alpha_\ell$ be an isomorphism $R_0\otimes\Z_\ell\xrightarrow{\sim}\Z_\ell^n$. The $\alpha_\ell$ correspond to
\[
\alpha:R_0\otimes\widehat{\Z}\xrightarrow{\sim}\widehat{\Z}^{n}.
\]
The identity embedding of $R_0$ into $\R^n$, together with $\alpha$, defines an isomorphism
\[
(1,\alpha):R_0\otimes(\R\times\widehat{\Z})\xrightarrow{\sim}(\R\times\widehat{\Z})^n,\quad\text{whence}
\]
\[
(1,\alpha)_{\Q}:(R_0\otimes\Q)\otimes\A\xrightarrow{\sim}\A^n,
\]
which makes $R=R_0\otimes\Q$ a lattice in $\A^n$.

\paragraph{Reminder 0.1.2.}{\itshape  The preceding construction is a bijection between the set of lattices $R_0\subset\R^n$ equipped with $\alpha:R_0\otimes\widehat{\Z}\xrightarrow{\sim}\widehat{\Z}^{n}$ and the set of lattices in $\A^n$. The inverse bijection associates to $R\subset\A^n$ the $\Z$-module $R_0=R\cap(\R\times\widehat{\Z})^n$, regarded as a lattice in $\R^n$ by the first projection $R_0\hookrightarrow\R^n$, and equipped with the second projection $\alpha$.\par}

\clearpage
\markboth{Authority 6 / printed 60}{}
\hypertarget{authority-6}{}
(0.1.3) Under this dictionary, the action of $\GL(n,\A)$ on the lattices in $\A^n$ is described as follows. Let $R$ be a lattice in $\A^n$ corresponding to $(R_0,\alpha)$, and let $g=(g_\infty,g_f)\in\GL(n,\A)=\GL(n,\R)\times\GL(n,\widehat{\Z}\otimes\Q)$.

Let $R'_0$ be the lattice isogenous to $R_0$ such that
\[
\alpha(R'_0\otimes\widehat{\Z})=g_f^{-1}(\widehat{\Z}^{n})\quad\text{in }\widehat{\Z}^{n}\otimes\Q.
\]
Then
\[
(g_\infty,g_f)(R_0,\alpha)=(g_\infty R'_0,g_f\alpha\circ g_\infty^{-1})
\]
and in particular
\[
\begin{aligned}
\text{for }g_\infty\in\GL(n,\R),\quad &(g_\infty,1)(R_0,\alpha)=(g_\infty R_0,\alpha\circ g_\infty^{-1})\\
\text{for }k\in\GL(n,\widehat{\Z}),\quad &(1,k)(R_0,\alpha)=(R_0,k\alpha).
\end{aligned}
\]
(0.1.4) Let $\mathcal G_{\A}\subset\Hom(\Q^n,\A^n)$ be the set of homomorphisms $g$ whose image is a lattice. Every $g$ extends by linearity to an automorphism $g_{\A}$ of $\A^n$, hence gives an isomorphism $\mathcal G_{\A}\sim\GL(n,\A)$. The set $\calR_{\A}$ of lattices in $\A^n$ is identified with the quotient
\[
\calR_{\A}=\GL(n,\A)/\GL(n,\Q).
\]
Likewise, the set $\calR_{\Z}$ of lattices in $\R^n$ is identified with the quotient
\[
\calR_{\Z}=\GL(n,\R)/\GL(n,\Z).
\]
Assertion 0.1.2 also says that the map
\[
\tag{0.1.4.1}
[\GL(n,\R)\times\GL(n,\widehat{\Z})]/\GL(n,\Z)\xrightarrow{\sim}\GL(n,\A)/\GL(n,\Q)
\]
is an isomorphism. To each adelic lattice, identified with $(R_0,\alpha)$, associate the lattice $R_0$. The resulting map from $\calR_{\A}$ to $\calR_{\Z}$ is identified with the composite
\[
\begin{array}{ccc}
\GL(n,\A)/\GL(n,\Q)&&\\
\uparrow\scriptstyle\sim&&\\
{}[\GL(n,\R)\times\GL(n,\widehat{\Z})]/\GL(n,\Z)&\longrightarrow&\GL(n,\R)/\GL(n,\Z)
\end{array}
\]

\clearpage
\markboth{Authority 7 / printed 61}{}
\hypertarget{authority-7}{}
It induces an isomorphism
\[
\tag{0.1.4.2}
\GL(n,\widehat{\Z})\backslash\calR_{\A}\iso\calR_{\Z}.
\]

\subsection*{0.2. Admissible representations}

\paragraph{0.2.1.} Let $F$ be a nonarchimedean local field and let $k$ be a field of characteristic $0$. An admissible representation of $\GL(2,F)$ over $k$ is a linear representation, finite- or infinite-dimensional,
\[
\pi:\GL(2,F)\longrightarrow\GL_k(V)
\]
such that the stabilizer of $v\in V$ is open, equivalently $\pi(g)v$ is continuous for the discrete topology of $V$, and such that, for $H\subset\GL(2,F)$ an open subgroup, the subspace $V^H$ of $H$-invariants in $V$ is finite-dimensional.

For the definition of the admissible contragredient $\check\pi$, see [9], p. 28.

\paragraph{0.2.2.} Rather than $\GL(2,\R)$, we consider the isomorphic group $\GL_{\R}(\C)$. This group contains $\C^*$, the maps $z\mapsto\lambda z$, hence $U_1\subset\C^*$, and complex conjugation $s$. Let
\[
K=U_1\cup sU_1
\]
be the normalizer of $U_1$ in $\SL_{\R}(\C)$. An admissible complex representation of $\GL_{\R}(\C)$ consists of:
\begin{enumerate}[label=(\alph*)]
\item a linear representation
\[
\pi|_K:K\longrightarrow\GL_{\C}(V),
\]
an algebraic direct sum of usual finite-dimensional irreducible representations of $K$, each occurring only finitely many times;
\item an action on $V$ of the Lie algebra $\mathfrak{gl}_{\R}(\C)$ of $\GL_{\R}(\C)$, extending the action of $\Lie(K)$ and such that, for $k\in K$ and $X\in\mathfrak{gl}_{\R}(\C)$,
\[
\tag{0.2.2.1}
\pi(k)\pi(X)\pi(k)^{-1}=\pi(\operatorname{ad}k.X).
\]
\end{enumerate}
It is easy to check that the given representation of $K$ extends uniquely to a representation of
\[
(\C^*\cup s\C^*)\subset\GL_{\R}(\C)
\]
which still satisfies (b).

\clearpage
\markboth{Authority 8 / printed 62}{}
\hypertarget{authority-8}{}
It is moreover enough to verify (0.2.2.1) for $k=s$.

\paragraph{0.2.3.} For the definition of admissible representations of $\GL(2,\A)$, we refer to [9], §9. Every irreducible admissible representation $\pi$ of $\GL(2,\A)$ on a complex vector space $V$ is a restricted tensor product: for every place $v$, there exists an irreducible admissible representation
\[
\pi_v:\GL(2,\Q_v)\longrightarrow\GL(V_v),
\]
and, for almost every prime number $p$, a vector $v_p\in V_p$, invariant under $\GL(2,\Z_p)$, such that
\[
V=\mathop{\bigotimes}\nolimits'_v V_v=\lim_{\longrightarrow,S}\bigotimes_{v\in S}V_v,
\]
the limit being over sufficiently large finite sets of places, with inclusions defined by the $v_p$. This is easy: see [9], 9.1. The same holds for $\GL(2,\A^f)$.

\section*{1. Lattice functions}

\subsection*{1.1. Functions on $\GL(2,\R)$}

\paragraph{1.1.1.} Let $(e_1,e_2)$ be the standard basis of $\Z^2$. The set of homomorphisms
\[
g:\Z^2\longrightarrow\C
\]
is identified with the set of pairs $(\omega_1,\omega_2)$ of complex numbers by
\[
g\longmapsto(g(e_1),g(e_2)).
\]
Let $G$ be the set of $g\in\Hom(\Z^2,\C)$ such that $g(\Z^2)$ is a lattice. One can identify $G$ with the set of pairs $(\omega_1,\omega_2)$ of complex numbers that are linearly independent over $\R$.

Every $g\in G$ extends linearly to $\R$- and $\C$-linear maps $g_{\R}$ and $g_{\C}$:
\[
\begin{gathered}
g_{\R}:\R^2\longrightarrow\C\\
g_{\C}:\C^2\longrightarrow\C.
\end{gathered}
\]
This gives various descriptions of $G$, which is identified respectively with:

\clearpage
\markboth{Authority 9 / printed 63}{}
\hypertarget{authority-9}{}
\begin{enumerate}[label=(\alph*)]
\item the set of lattices in $\C$ endowed with a basis, by
\[
g\longmapsto (g(\Z^2);\text{ basis }(g(e_1),g(e_2)));
\]
\item $\Isom_{\R}(\R^2,\C)$, by $g\mapsto g_{\R}$;
\item[(b')] the set of pairs consisting of a complex structure on $\R^2$, and of a $\C$-linear isomorphism, for this complex structure, from $\R^2$ to $\C$;
\item[(c)] the open subset of $\Hom_{\C}(\C^2,\C)$ formed by the forms $\ell$ whose restriction to $\R^2$ is injective, by $g\mapsto g_{\C}$.
\end{enumerate}

The interpretations (a) and (c) display a complex structure on $G$ for which the functions
\[
\omega_i=g(e_i)
\]
are holomorphic.

Interpretation (b) displays a right action, by composition, of $\GL(2,\R)$ on $G$. This action respects the complex structure of $G$ and, by (a),
\[
G/\GL(2,\Z)
\]
is identified with the set of lattices in $\C$, endowed with its natural complex structure. In the coordinate system $(\omega_1,\omega_2)$, the action is written
\[
(\omega_1,\omega_2),
\begin{pmatrix}a&b\\ c&d\end{pmatrix}
\longmapsto
(\omega_1,\omega_2)\begin{pmatrix}a&b\\ c&d\end{pmatrix},
\]
with matrix multiplication.

Interpretation (b) also displays a left action of $\GL_{\R}(\C)$ on $G$:
\[
g,\ (\omega_1,\omega_2)\longmapsto(g(\omega_1),g(\omega_2)).
\]
Identify $\C^*$ with a subgroup of $\GL_{\R}(\C)$ by
\[
\lambda\longmapsto(z\longmapsto\lambda z).
\]
The action of $\C^*\subset\GL_{\R}(\C)$ on $G$ is holomorphic, and $G$ is a principal homogeneous space with group $\C^*$ over
\[
X^\pm=\C^*\backslash G.
\]
By (b'), $X^\pm$ is identified with the set of complex structures on $\R^2$. The quotient complex structure induced from that of $G$ on $X^\pm$ is also deduced from the following inclusion
\[
X^\pm\hookrightarrow\mathbb P^1(\C):
\]
to $x\in X^\pm$, one associates the kernel of the $\C$-linear map

\clearpage
\markboth{Authority 10 / printed 64}{}
\hypertarget{authority-10}{}
\[
\C^2=\R^2\otimes\C\longrightarrow (\R^2,x)
\]
extending the identity
\[
\R^2\longrightarrow\R^2.
\]

\paragraph{1.1.2.} Let $g_0\in G$ be the lattice with basis $(i,1)$. We shall choose $g_0$ as the origin in $G$, and identify $G$ with $\GL(2,\R)$ by
\[
h\in\GL(2,\R)\longmapsto g_0h.
\]
The right or left actions of $\GL(2,\R)$ or $\GL_{\R}(\C)$ introduced above are then identified with right or left translations, $\GL_{\R}(\C)$ being identified with $\GL(2,\R)$ by means of
\[
(i,1):\R^2\iso\C.
\]

\paragraph{1.1.3.} On $\GL(2,\R)$, define a function $\|g\|$ by
\[
\|g\|=\Tr({}^{t}gg+({}^{t}gg)^{-1})
=(a^2+b^2+c^2+d^2)(1+\det(g)^{-2}).
\]
A function $f$ on $\GL(2,\R)$ will be said to be $C^\infty$ with \emph{moderate growth} if there exist $A>0$, $N>0$, such that
\[
\tag{1.1.3.1}
f(g)\le A\|g\|^N,
\]
and all derivatives of $f$, relative to invariant vector fields, satisfy analogous conditions. Condition (1.1.3.1) amounts to saying that $f$ is bounded above by an algebraic function. This is a very weak growth condition (the analogue for $\GL(1,\R)$ is: bounded above by $A.\chi$, for a suitable quasi-character $\chi$).

By (1.1.2), this terminology is transported to functions on $G$, and it is extended to functions on the set $G/\GL(2,\Z)$ of lattices in $\C$, identified with certain functions on $G$.

\paragraph{Definition 1.1.4.}{\itshape  A holomorphic modular form of weight $k$, of group $\GL(2,\Z)$, is a lattice function $f(R)$
\[
f:G/\GL(2,\Z)\longrightarrow\C
\]
which is\par}

\clearpage
\markboth{Authority 11 / printed 65}{}
\hypertarget{authority-11}{}
{\itshape \begin{enumerate}[label=(\alph*)]
\item holomorphic,
\item such that
\[
f(\lambda R)=\lambda^{-k}f(R)\qquad(R\text{ is a lattice and }\lambda\in\C^*),
\]
\item[(d)] of moderate growth.
\end{enumerate}
\par}
When \(f\) is identified with a function on \(G\), these conditions become:
\[
\begin{array}{ll}
(a_1)& f\text{ is holomorphic},\\[0.25em]
(b_1)& f(\lambda g)=\lambda^{-k}f(g)\qquad(g\in G,\ \lambda\in\C^*),\\[0.25em]
(c_1)& f(g\gamma)=f(g)\qquad(g\in G,\ \gamma\in\GL(2,\Z)),\\[0.25em]
(d_1)& f\text{ is of moderate growth.}
\end{array}
\]

\paragraph{1.1.5.} Two coordinate systems on \(G\) are very convenient. In each of them it is useful to know:
\begin{enumerate}[label=(\Alph*)]
\item the invariant volume element, for \(\GL(2,\R)\) on the right and \(\GL_{\R}(\C)\) on the left;
\item the value of \(\det g\), to which a meaning is given by 1.1.2;
\item the left action of \(\C^*\), and the right action of \(\GL(2,\R)\);
\item the Casimir operator, invariant under left and right translations. It is the central element
\[
\Omega=\frac12H^2+XY+YX
\]
of the enveloping algebra of the Lie algebra of \(\GL(2,\R)\), acting by convolution.
\end{enumerate}
We have put, in the complexification of this Lie algebra,
\[
H=\begin{pmatrix}0&-i\\ i&0\end{pmatrix},\qquad
X=\frac12\begin{pmatrix}-1&-i\\ -i&1\end{pmatrix},\qquad
Y=\frac12\begin{pmatrix}-1&i\\ i&1\end{pmatrix}.
\]
In the basis
\[
\frac12\begin{pmatrix}-i\\1\end{pmatrix},\qquad \frac12\begin{pmatrix}i\\1\end{pmatrix},
\]
identified with the basis dual to the basis \(z,\overline z\) of \(\Hom_{\R}(\C,\C)\), one has
\[
H=\begin{pmatrix}1&0\\0&-1\end{pmatrix},\qquad
X=\begin{pmatrix}0&1\\0&0\end{pmatrix},\qquad
Y=\begin{pmatrix}0&0\\1&0\end{pmatrix}.
\]
\(H\) is the infinitesimal generator of the compact torus \(U_1\subset\C^*\) of \(\GL_{\R}(\C)\simeq\GL(2,\R)\).

\clearpage
\markboth{Authority 12 / printed 66}{}
\hypertarget{authority-12}{}
\noindent(isomorphism 1.1.2).

\paragraph{1.1.5.1. Homogeneous coordinates.}
\[
g\longmapsto(\omega_1,\omega_2)=(g(e_1),g(e_2)).
\]
Put
\[
\omega_1=x_1+iy_1,\qquad \omega_2=x_2+iy_2.
\]
For \(\beta\) a holomorphic form, write \(|\beta|^2\) for the real form
\[
\beta\wedge\overline\beta
\]
and for the corresponding positive volume element when \(\beta\) is a \(2\)-form.

\[
(A_h)\qquad
dg=\left|\frac{d\omega_1\wedge d\omega_2}{\omega_1\wedge\omega_2}\right|^2
=
4\,dx_1\,dy_1\,dx_2\,dy_2\,(x_1y_2-x_2y_1)^{-2}.
\]
\[
(B_h)\qquad
\det g=x_2y_1-x_1y_2.
\]
\[
(C_h)\qquad
\lambda\cdot(\omega_1,\omega_2)=(\lambda\omega_1,\lambda\omega_2),
\]
\[
(\omega_1,\omega_2)\gamma=(a\omega_1+b\omega_2,\ c\omega_1+d\omega_2)
\quad\text{for}\quad
\gamma=\begin{pmatrix}a&c\\ b&d\end{pmatrix}.
\]
\[
(D_h)\qquad
H*f=\left[\left(-\omega_1\frac{\partial}{\partial\omega_1}
+\overline\omega_1\frac{\partial}{\partial\overline\omega_1}\right)
+\left(-\omega_2\frac{\partial}{\partial\omega_2}
+\overline\omega_2\frac{\partial}{\partial\overline\omega_2}\right)\right]f,
\]
\[
X*f=\left[-\overline\omega_1\frac{\partial}{\partial\omega_1}
-\overline\omega_2\frac{\partial}{\partial\omega_2}\right]f,
\]
\[
Y*f=\left[-\omega_1\frac{\partial}{\partial\overline\omega_1}
-\omega_2\frac{\partial}{\partial\overline\omega_2}\right]f.
\]

\paragraph{1.1.5.2. Poincaré half-planes.}
\[
g\longmapsto[\lambda,z]=[\omega_2,\omega_1/\omega_2],\qquad
(\omega_1,\omega_2)=\lambda(z,1).
\]
Put \(z=x+iy\). In the coordinate system \([\lambda,z]\), \(\lambda\) runs through \(\C^*\), and \(z\) runs through the set of complex numbers with imaginary part \(y\ne0\).

\clearpage
\markboth{Authority 13 / printed 67}{}
\hypertarget{authority-13}{}
\[
(A_p)\qquad
dg=\left|\frac{d\lambda}{\lambda}\right|^2\,y^{-2}|dz|^2
=
2\,\frac{d\lambda}{\lambda}\frac{d\overline\lambda}{\overline\lambda}\,y^{-2}\,dx\,dy.
\]
\[
(B_p)\qquad
\det g=|\lambda|^2y.
\]
\[
(C_p)\qquad
\lambda[\mu,z]=[\lambda\mu,z],
\]
\[
[\lambda,z]\gamma=
\left[\lambda(cz+d),\,\frac{az+b}{cz+d}\right]
\quad\text{for}\quad
\gamma=\begin{pmatrix}a&c\\ b&d\end{pmatrix}.
\]
\[
(D_p)\qquad
H*f=\left[-\lambda\frac{\partial}{\partial\lambda}
+\overline\lambda\frac{\partial}{\partial\overline\lambda}\right]f,
\]
\[
X*f=\left[2y\,\frac{\overline\lambda}{\lambda}\frac{\partial}{\partial z}
-\overline\lambda\frac{\partial}{\partial\lambda}\right]f,
\]
\[
Y*f=\left[-2y\,\frac{\lambda}{\overline\lambda}\frac{\partial}{\partial\overline z}
-\lambda\frac{\partial}{\partial\overline\lambda}\right]f,
\]
\[
\Omega*f=
\left[
\frac12\left(\lambda\frac{\partial}{\partial\lambda}
+\overline\lambda\frac{\partial}{\partial\overline\lambda}\right)^2
+\left(\lambda\frac{\partial}{\partial\lambda}
+\overline\lambda\frac{\partial}{\partial\overline\lambda}\right)
+4\left(y\lambda\frac{\partial}{\partial z}\frac{\partial}{\partial\lambda}
-y\overline\lambda\frac{\partial}{\partial\overline z}\frac{\partial}{\partial\overline\lambda}\right)
-8y^2\frac{\partial}{\partial z}\frac{\partial}{\partial\overline z}
\right]f.
\]

\paragraph{1.1.6.} Formulas (1.1.5.2) make clear that \(G\) is a principal homogeneous space with group \(\C^*\) over the double Poincaré half-plane
\[
X^\pm=\{z\in\C\mid \operatorname{Im}(z)\ne0\}.
\]
Restriction to the section \(z\mapsto(z,1)\) of \(G\to X^\pm\) identifies the set \(\mathfrak F_k\) of functions satisfying \((b_1)\) (1.1.4) with the set of functions on
\[
X^\pm=\C-\R.
\]

\clearpage
\markboth{Authority 14 / printed 68}{}
\hypertarget{authority-14}{}
Under this identification, holomorphic functions correspond to holomorphic functions. The right action of \(\GL(2,\R)\) on \(G\) induces a left action on functions on \(G\), by
\[
(\gamma f)(g)=f(g\gamma).
\]
This action preserves \(\mathfrak F_k\). When \(\mathfrak F_k\) is identified with functions on \(X^\pm\), this action becomes, for
\[
\gamma=\begin{pmatrix}a&c\\ b&d\end{pmatrix},
\]
\[
(\gamma f)(z)=(cz+d)^{-k}\,
f\!\left(\frac{az+b}{cz+d}\right).
\]
It is traditional, in order to have a right action, to put
\[
f|\gamma={}^{t}\gamma(f)
\quad
\left(\text{for }\gamma=\begin{pmatrix}a&b\\ c&d\end{pmatrix},\ 
{}^{t}\gamma=\begin{pmatrix}a&c\\ b&d\end{pmatrix}\right).
\]

Holomorphic modular forms of weight \(n\), of group \(\GL_2(\Z)\), are therefore identified with functions \(f\) on \(X^\pm\) such that:
\[
(a_2)\quad f\text{ is holomorphic},
\]
\[
(c_2)\quad
\text{for }\gamma=\begin{pmatrix}a&b\\ c&d\end{pmatrix}\in\GL(2,\Z),\qquad
f\!\left(\frac{az+b}{cz+d}\right)=(cz+d)^k f(z),
\]
\[
(d_2)\quad
|f(x+iy)|\le Ay^N
\quad\text{for }\Re(x)\text{ bounded},\ y\to\infty,\ N\text{ suitable}.
\]
It follows from reduction theory that this condition implies the apparently stronger condition \((d_1)\).

\paragraph{Variant 1.1.7.} The elements of \(\GL(2,\Z)\) of determinant \(-1\) permute the two connected components of \(X^\pm\). Thus the forms under consideration may also be identified with functions on the Poincaré half-plane
\[
X=\{z\in\C\mid \operatorname{Im}(z)>0\},
\]
satisfying \((a_2)\), \((d_2)\), and \((c_2)\) for \(\gamma\in\SL(2,\Z)\).

\paragraph{Variant 1.1.8.} Let \(SG\) be the set of \(g\in G\) such that
\[
g(e_1)\wedge g(e_2)=i\wedge1
\]
in \(\bigwedge^2_{\R}\C\). On \(SG\subset G\), the subgroup \(\SL(2,\R)\) of \(\GL(2,\R)\) acts on the right, and \(\SL_{\R}(\C)\supset U_1\) on the left. A modular form as above is determined by its restriction to \(SG\).

\clearpage
\markboth{Authority 15 / printed 69}{}
\hypertarget{authority-15}{}
Let \(\widetilde{SG}\) be the set of pairs consisting of \(g\in SG\) and an isomorphism, over \(g_{\R}\), of the universal covering of \(\R^2-\{0\}\), with base point
\[
e_2=(0,1),
\]
with that of \(\C-\{0\}\), with base point \(1\). \(\widetilde{SG}\) is a universal covering of \(SG\). The universal covering \(\widetilde{\SL}(2,\R)\) of \(\SL(2,\R)\) acts on the right, and that of \(\SL_{\R}(\C)\) on the left. In particular, the universal covering
\[
\widetilde U_1\subset\widetilde{\SL}_{\R}(\C)
\]
acts on the left. Note that, for \(k\in\C\) and \(\lambda\) in \(\widetilde{\C^*}\), in particular in \(\widetilde U_1\), \(\lambda^k\) is well-defined.

For \(k\in\R\) and \(\Gamma\) a discrete subgroup of \(\widetilde{\SL}_2(\R)\), a holomorphic modular form of weight \(k\), of group \(\Gamma\), is a function \(f\) on \(\widetilde{SG}\) such that:
\[
(a_3)\quad
f\text{ satisfies a holomorphy condition;}
\]
more precisely, locally, when \(f\) is brought back to a function on \(SG\) by a local section of \(\widetilde{SG}\to SG\), and when it is extended to \(G\) as \(\widetilde f\), so that
\[
\widetilde f(\lambda g)=\lambda^{-k}f(g)\qquad(\lambda\in\R_+^*),
\]
one wants \(\widetilde f\) to be holomorphic;
\[
(b_3)\quad f(\lambda g)=\lambda^{-k}f(g)\qquad(\lambda\in\widetilde U_1),
\]
\[
(c_3)\quad f(g\gamma)=f(g)\qquad(\gamma\in\Gamma),
\]
\[
(d_3)\quad |f(g)|,\text{ which by }(b_3)\text{ comes from a function on }SG,\text{ is of moderate growth.}
\]

\section*{1.2. Functions on \(\GL(2,\A)\)}

\paragraph{1.2.1.} Let
\[
\GA\subset\Hom(\Q^2,\C\times(\Af)^2)
\]
be the set of homomorphisms \(g\) whose image is a lattice; see (0.1.1). Here \(\C\times(\Af)^2\) is a free \(\A\)-module of rank two. These \(g\) are identified with isomorphisms
\[
\A^2\iso\C\times(\Af)^2,
\]
whence a right action of \(\GL(2,\A)\) and a left action of
\[
\GL_{\R}(\C)\times\GL(2,\Af).
\]
The quotient
\[
\GA/\GL(2,\Q)
\]
is the set of lattices in \(\C\times(\Af)^2\). By (0.1.2), the set of lattices in \(\C\times(\Af)^2\) is identified with the set of pairs \((R_0,\alpha)\), consisting of a lattice \(R_0\subset\C\) and of
\[
R_0\otimes\widehat{\Z}\iso\widehat{\Z}^2.
\]
A lattice function \(f(R_0)\), \(R_0\subset\C\), therefore defines an adelic lattice function
\[
R\longmapsto(R_0,\alpha)\longmapsto f(R_0)
\qquad(R\subset\C\times(\Af)^2).
\]
This construction

\clearpage
\markboth{Authority 16 / printed 70}{}
\hypertarget{authority-16}{}
identifies holomorphic modular forms of weight $k$ and group $\GL(2,\Z)$ (1.1.4) with functions on $\GA$ which are:
\[
(a_4)\quad \text{holomorphic in }g_\infty,
\]
\[
({}_{\infty}b_4)\quad f(\lambda g)=\lambda^{-k}f(g)\qquad(\lambda\in\C^*\subset\GL_{\R}(\C)),
\]
\[
({}_{f}b_4)\quad f(\varkappa g)=f(g)\qquad(\varkappa\in\GL(2,\widehat{\Z})\subset\GL(2,\Af)),
\]
\[
(c_4)\quad f(g\gamma)=f(g)\qquad(\gamma\in\GL(2,\Q)),
\]
\[
(d_4)\quad \text{of moderate growth as a function of }g_\infty.
\]

\paragraph{Definition 1.2.2.}{\itshape  A holomorphic modular form of weight $k$ is a function on $\GA$ satisfying conditions $(a_4)$ through $(d_4)$ above, with $({}_{f}b_4)$ replaced by
\[
({}_{f}b_5)\quad \text{for }\varkappa\text{ in a suitable compact open subgroup }K\text{ of }\GL(2,\Af),\text{ one has } f(\varkappa g)=f(g).
\]
\par}
We shall say that $f$ is of \emph{level} $K$ if $f$ is invariant under $K$.

One again checks that it is enough to verify $(d_4)$ in a Siegel domain.

\paragraph{1.2.3.} Let us translate this definition into more concrete terms, for $K\subset\GL(2,\widehat{\Z})$. There is then an integer $n$ and $H\subset\GL(2,\Z/n)$ such that $K$ is the inverse image of $H$ under the surjective map
\[
\begin{tikzcd}[row sep=large,column sep=large]
\GL(2,\widehat{\Z}) \arrow[r] & \GL(2,\Z/n)\\
K \arrow[u,hook] \arrow[r] & H \arrow[u,hook]
\end{tikzcd}
\]
Let $K_n$ be defined by $H=\{e\}$. By (0.1.4.1), the quotient
\[
K_n\backslash\GA/\GL(2,\Q)
\]
is identified with the set of pairs $(R_0,\bar\alpha)$ consisting of a lattice $R_0\subset\C$ and of
\[
\bar\alpha:R_0/nR_0\iso(\Z/n\Z)^2.
\]
A holomorphic modular form of weight $k$ will be said to be of \emph{level} $n$ if it is invariant under $K_n$. Such a form is identified with

\clearpage
\markboth{Authority 17 / printed 71}{}
\hypertarget{authority-17}{}
a holomorphic function of marked lattices
\[
f(R_0,\bar\alpha),
\]
homogeneous of weight $-k$, and satisfying a suitable growth condition.

More generally, for $K$ defined by $(N,H)$, the quotient
\[
K\backslash\GA/\GL(2,\Q)
\]
is identified with the set of pairs $(R_0,\bar\alpha)$ consisting of a lattice $R_0\subset\C$ and a coset
\[
\bar\alpha\in H\backslash\Isom(R_0/nR_0,(\Z/n\Z)^2).
\]

\paragraph{1.2.4.} Note that the set of marked lattices $K\backslash\GA/\GL(2,\Q)$ is in general disconnected. By what precedes, or by (0.1.4.1), it is indeed identified with
\[
\tag{1.2.4.1}
\begin{tikzcd}[row sep=large,column sep=small]
K\backslash\GA/\GL(2,\Q) & K\backslash\bigl(G\times\GL(2,\widehat{\Z})\bigr)/\GL(2,\Z) \arrow[l,"\sim"'] \arrow[d,"\sim"] \arrow[r,"\sim"] & {}\\
& G\times\bigl(H\backslash\GL(2,\Z/n)\bigr)/\GL(2,\Z) &
\end{tikzcd}
\]
The two connected components of $G$ are interchanged by the elements of $\GL(2,\Z)$ of determinant $-1$, and $\SL(2,\Z)$ maps to $\SL(2,\Z/n)$. From (1.2.4.1) one therefore obtains a bijection
\[
\pi_0\bigl(K\backslash\GA/\GL(2,\Q)\bigr)
\xrightarrow[\sim]{\mu}
\det(H)\backslash(\Z/n\Z)^*.
\]
or, in adelic terms,
\[
\pi_0\bigl(K\backslash\GA/\GL(2,\Q)\bigr)
\iso
\det K\backslash\{\pm1\}\times(\Af)^*/\Q^*.
\]
In terms of marked lattices $(R_0,\bar\alpha)$, this map is described as follows. The lattice $R_0\subset\C$ has a natural orientation: the generator $e$ of $\bigwedge^2R_0$ whose image in $\bigwedge^2_{\R}\C$ is a positive multiple of $i\wedge1$. If
\[
\alpha:R_0/nR_0\iso(\Z/n\Z)^2
\]
is a representative of $\bar\alpha$, then $\det(\alpha)=\bigwedge^2\alpha(e)$ is an element of $(\Z/n\Z)^*$ whose class $\mu(R_0,\alpha)$ modulo $\det(H)$ depends only on $\bar\alpha$.

\clearpage
\markboth{Authority 18 / printed 72}{}
\hypertarget{authority-18}{}
\paragraph{1.2.5.} Let $f$ be a holomorphic modular form of weight $k$ and level $K$, with $K$ as in 1.2.3. By (1.2.4.1), $f$ is determined by the family of functions
\[
f(z;\sigma)\qquad(z\in X^\pm,
\ \sigma\in\GL(2,\Z/n)),
\]
defined by
\[
f(z;\sigma)=f\bigl((z,1),\sigma\bigr).
\]
The variance properties of $f$ are written
\[
\tag{1.2.5.1}
f(z;\sigma)=f(z;h\sigma)\qquad(h\in H),
\]
\[
\tag{1.2.5.2}
f(z;\sigma)=(cz+d)^{-k}
 f\left(\frac{az+b}{cz+d},\sigma\cdot\widetilde\gamma\right)
\]
for
\[
\gamma=\begin{pmatrix}a&c\\ b&d\end{pmatrix}\in\GL(2,\Z),
\]
and $\widetilde\gamma$ its reduction modulo $n$.

The form $f$ is also determined by the family of restrictions to the Poincaré half-plane $X$ of the $f(z,\sigma)$, with (1.2.5.2) now needing to be verified only for $\gamma\in\SL(2,\Z)$.

Let $\sigma\in\GL(2,\Z/n)$, and let $\Gamma_\sigma\subset\SL(2,\Z)$ be the inverse image of $\sigma^{-1}H\sigma$. Conditions (1.2.5.1) and (1.2.5.2) imply
\[
\tag{1.2.5.3}
f(z;\sigma)=(cz+d)^{-k}f\left(\frac{az+b}{cz+d},\sigma\right)
\qquad
\left(\begin{pmatrix}a&c\\ b&d\end{pmatrix}\in\Gamma_\sigma\right).
\]
Conversely, if $\Sigma$ is a subset of $\GL(2,\Z/n)$ mapping bijectively onto $\det H\backslash(\Z/n\Z)^*$, then $f$ is uniquely determined by the $f(z,\sigma)$, with $z\in X$ and $\sigma\in\Sigma$, subject to (1.2.5.3). In particular, if
\[
\det:H\longrightarrow(\Z/n\Z)^*
\]
is surjective, then $f$ is uniquely determined by
\[
f(z)=f(z;1)\qquad(z\in X),
\]
subject to (1.2.5.3), for $\sigma=1$.

\paragraph{Variant.} Let
\[
\SG_{\A}=\SG\times\SL(2,\Af).
\]
Let $\widetilde{\SG}_2$ be the double cover of $\SG$, quotient of the universal cover (1.1.8), and let $\widetilde{\SL}(2,\Q_p)$ be the double cover of $\SL(2,\Q_p)$, a central extension by $\Z/2$. It is known that, for almost all $p$, $\widetilde{\SL}(2,\Q_p)$ splits above $\SL(2,\Z_p)$. A restricted product
\[
\SG'_{\A}=\widetilde{\SG}_2\times\prod_p\SL(2,\Q_p)
\]
is therefore defined; it is a group cover by $\sum_v(\Z/2)$ of $\SG_\A$. Let $\widetilde{\SG}_{\A}$ be the double cover deduced from it via the summation map
\[
\sum_v\Z/2\longrightarrow\Z/2.
\]
A double cover of $\SL(2,\A)$ acts on the right on $\widetilde{\SG}_\A$.

It is

\clearpage
\markboth{Authority 19 / printed 73}{}
\hypertarget{authority-19}{}
known (C. C. Moore, Publ. Math. IHES 35) that this cover splits over $\SL(2,\Q)$, so that a right action of $\SL(2,\Q)$ is defined. Similarly, a double cover of $\SL_{\R}(\C)\times\SL(2,\Af)$ acts on the left. For $k\in\tfrac12\Z$ and $K$ a compact open subgroup of the induced double cover of $\SL(2,\Af)$, a holomorphic modular form of weight $k$ and level $K$ is a function on $\widetilde{\SG}_\A$:
\[
(a_6)\quad\text{satisfying a holomorphy condition},
\]
\[
({}_{\infty}b_6)\quad f(\lambda g)=\lambda^{-k}f(g)\qquad(\lambda\in\widetilde U_1),
\]
\[
({}_{f}b_6)\quad f(\varkappa g)=f(g)\qquad(\varkappa\in K),
\]
\[
(c_6)\quad f(g\gamma)=f(g)\qquad(\gamma\in\SL(2,\Q)),
\]
\[
(d_6)\quad\text{of moderate growth as a function of }g_\infty.
\]
The quotient $\widetilde{\SG}_{\A}/\SL(2,\Q)$ is connected; these forms are identified with holomorphic functions on $X$.

As in 1.1.2, identify $\widetilde{\SG}_{\A}$ and the preceding double cover of $\SL_{\R}(\C)\times\SL(2,\Af)$ with a double cover $\widetilde{\SL}(2,\A)$ of $\SL(2,\A)$. It is possible to make a group larger than $\widetilde{\SL}(2,\A)$ act on the space $L$ of functions on $\widetilde{\SG}_{\A}/\SL(2,\Q)$: $\GL(2,\Q)$ acts by conjugation on $\SL_2$, hence by transport of structure on $\widetilde{\SL}(2,\A)$, and this makes it possible to define a group
\[
\GL(2,\Q)\cdot\widetilde{\SL}(2,\A)
\]
containing $\GL(2,\Q)$ and $\widetilde{\SL}(2,\A)$, the intersection of the two being $\SL(2,\Q)$. One then has
\[
\widetilde{\SL}(2,\A)/\SL(2,\Q)
\iso
\GL(2,\Q)\cdot\widetilde{\SL}(2,\A)/\GL(2,\Q),
\]
so that $\GL(2,\Q)\cdot\widetilde{\SL}(2,\A)$ acts on $L$, in fact through
\[
\GL(2,\Q)\cdot\widetilde{\SL}(2,\A)/\Q^*.
\]

\clearpage
\markboth{Authority 20 / printed 74}{}
\hypertarget{authority-20}{}
\section*{1.3. Cusps}

\paragraph{1.3.1.} Let $f$ be a holomorphic modular form of weight $k$ and level $n$, described as in 1.2.5 by functions
\[
f(z;\sigma)\qquad(z\in X,\ \sigma\in\GL(2,\Z/n)).
\]
By (1.2.5.3), one has
\[
f(z+n;\sigma)=f(z;\sigma).
\]
Put
\[
q^\lambda=e^{2\pi i\lambda z}.
\]
The functions $f(z;\sigma)$ then have Fourier expansions
\[
\tag{1.3.1.1}
f(z;\sigma)=\sum_{\lambda}a(\lambda,f,\sigma)q^\lambda\qquad(\lambda\in\tfrac1n\Z).
\]
One immediately checks that the growth condition $(d_4)$ is equivalent to
\[
\tag{1.3.1.2}
a(\lambda,f,\sigma)=0\qquad\text{for }\lambda<0.
\]
The form $f$ is called \emph{cuspidal} if the $a(0,f,\sigma)$ vanish, i.e. if the $f(z,\sigma)$ are “zero at the cusps.” One has
\[
a(0,f,\sigma)=\int_{\R/n\Z}f(z+u)\,\frac{du}{n},
\]
whence the following adelic writing.

\paragraph{Reminder 1.3.2.}{\itshape  Let $U$ be the lower unipotent subgroup
\[
\begin{pmatrix}1&0\\ *&1\end{pmatrix}.
\]
The form $f$ is \emph{cuspidal} if and only if
\[
\tag{1.3.2.1}
f_P(g)=\int_{U_\A/U_\Q}f(gu)\,du=0.
\]
\par}
Let $P$ be the lower triangular subgroup. One has
\[
f_P(gp)=f_P(g)\qquad(p\in P(\Q)U(\A)),
\]
\[
f_P(\lambda g)=\lambda^{-k}f_P(g),
\]
\[
f_P(\varkappa g)=f_P(g)\qquad(\varkappa\in K_n).
\]
Consequently, to verify that $f_P(g)$ is zero, it is enough to verify it for
\[
g\in(\text{the image of }X\text{ in }G\text{ by }z\mapsto(z,1))\times\GL(2,\widehat{\Z}),
\]
and one obtains the condition written above.

\clearpage
\markboth{Authority 21 / printed 75}{}
\hypertarget{authority-21}{}
Of course, (1.3.2.1) will hold for every unipotent subgroup of \(\GL(2,\Q)\), since these are all conjugate.

\paragraph{(1.3.3)} Let \(f_1\) and \(f_2\) be two holomorphic modular forms of weight \(k\), with \(f_1\) cuspidal. Let \(n\) be such that they are of level \(n\). The function
\[
\Phi(g)=f_1(g)\overline{f_2(g)}
\]
on \(G_{\mathbb A}/\GL(2,\Q)\) satisfies
\[
\Phi\!\left(\begin{pmatrix}x&0\\0&x\end{pmatrix}g\right)=\|x\|^{2k}\Phi(g),
\]
so that
\[
\|\det(g)\|^{-k}f_1(g)\overline{f_2(g)}
\]
is invariant under the center \(Z_{\mathbb A}\). Put
\[
\langle f_1,f_2\rangle
=
\int_{Z_{\mathbb A}\backslash G_{\mathbb A}/\GL(2,\Q)}
\|\det(g)\|^{-k}f_1(g)\overline{f_2(g)}\,dg.
\]
This is the \emph{Petersson scalar product}, left-invariant under \(\GL(2,\mathbb A)\). By (1.1.5.2), up to a constant independent of \(k\) and \(n\), it is written
\[
\langle f_1,f_2\rangle
=
\frac1{\#\GL(2,\Z/n)}
\sum_{\sigma\in\Sigma}
\int_{X/\Gamma}
y^{-k-2}f_1(z;\sigma)\overline{f_2(z;\sigma)}\,dz\,d\bar z,
\]
where
\[
\Gamma=\Ker\bigl(\SL(2,\Z)\longrightarrow\SL(2,\Z/n)\bigr)
\]
and where \(\det\) maps \(\Sigma\subset\GL(2,\Z/n)\) bijectively onto \((\Z/n\Z)^*\). If \(D\) is a fundamental domain for \(\SL(2,\Z)\), one may also write
\[
\langle f_1,f_2\rangle
=
\frac1{\#\GL(2,\Z/n)}
\sum_{\sigma}
\int_D
y^{-k-2}f_1(z;\sigma)\overline{f_2(z;\sigma)}\,dz\,d\bar z,
\]
where \(\sigma\) runs through all of \(\GL(2,\Z/n)\). In this form, it is clear that the integral converges as soon as \(f_1\) is cuspidal: for \(D\) the standard domain, \(f_1\) decreases at least like \(e^{-y/n}\) on \(D\), as \(y\to\infty\), whereas \(f_2\) is bounded.

\clearpage
\markboth{Authority 22 / printed 76}{}
\hypertarget{authority-22}{}
\section*{2. Modular forms and the spectrum of \(\GL(2)\)}

\subsection*{2.1. Holomorphic modular forms and representations of \(\GL(2,\R)\)}

\paragraph{(2.1.1)} Let \(s\in \GL_{\R}(\C)\) be complex conjugation, and resume the notation of 1.1.5. For every integer \(k\ge1\), denote by \(D_{k-1}\) the following irreducible admissible representation of \(\GL_{\R}(\C)\):
\begin{enumerate}[label=(\alph*)]
\item \(D_{k-1}\) has a basis \(e_n\), indexed by the integers \(n\equiv k\pmod 2\) such that \(|n|\ge k\).
\item The actions of \(\mathfrak{gl}_{\R}(\C)\) and of \(\C^*\cup s\C^*\subset \GL_{\R}(\C)\) satisfy
\[
\tag{b1}
\lambda*e_n=\lambda^k(\lambda\bar\lambda^{-1})^{(n-k)/2}e_n
\qquad(\lambda\in\C^*),
\]
\[
\tag{b2}
s*e_n=e_{-n},
\]
\[
\tag{b3}
H*e_n=n e_n
\qquad\text{(follows from (b1)),}
\]
\[
\tag{b4}
X*e_n=\frac{k+n}{2}e_{n+2}
\qquad(0\text{ if }n=-k),
\]
\[
\tag{b5}
Y*e_n=\frac{k-n}{2}e_{n-2}
\qquad(0\text{ if }n=k).
\]
\end{enumerate}

Let \(f\) be a function on \(G\) satisfying
\[
f(\lambda g)=\lambda^{-k}f(g),
\]
i.e. \(\lambda*f=\lambda^k f\) for \(\lambda\in\C^*\subset \GL_{\R}(\C)\). From the formulas (1.1.5.2), one checks that \(f\) is holomorphic if and only if \(Y*f=0\). For \(f\) holomorphic, put
\[
e_n=\frac{(k-1)!}{((n+k-2)/2)!}X^{(n-k)/2}*f
\qquad(n\ge k,\ n\equiv k\pmod2),
\]
\[
e_{-n}=s*e_n.
\]
One has \(e_k=f\), and the \(e_n\) satisfy formulas (b). The function \(f\) therefore generates an irreducible admissible representation of \(\GL_{\R}(\C)\), of type \(D_{k-1}\). Up to a scalar, it is the unique vector of this representation satisfying \(Y*f=0\). A converse is checked easily.

Let \(\mathfrak F(G,\GL(2,\Z))\) be the set of \(C^\infty\) functions of moderate growth on \(G\), right-invariant under \(\GL(2,\Z)\).

\clearpage
\markboth{Authority 23 / printed 77}{}
\hypertarget{authority-23}{}
Definition 1.1.4 may be reformulated as follows.

\paragraph{Scholium 2.1.2.}{\itshape  The set of holomorphic modular forms of weight \(k\) and group \(\GL(2,\Z)\) is identified, by
\[
\varphi\longmapsto\varphi(e_k),
\]
with the set of homomorphisms of representations
\[
\Hom_{\GL_{\R}(\C)}\bigl(D_{k-1},\mathfrak F(G,\GL(2,\Z))\bigr).
\]
\par}
“Homomorphism of representations” means a linear map which commutes with the actions of \(K\) and \(\mathfrak{gl}_{\R}(\C)\).

In the adelic framework, let \(\mathfrak F(G_{\mathbb A})\) likewise be the set of functions on \(G_{\mathbb A}\), right-invariant under \(\GL(2,\Q)\), left-invariant under a sufficiently small open compact subgroup, and \(C^\infty\) of moderate growth in \(g_\infty\). Again one has:

\paragraph{Scholium 2.1.3.}{\itshape  The set of holomorphic modular forms of weight \(k\) (1.2.2) is identified with
\[
\Hom_{\GL_{\R}(\C)}\bigl(D_{k-1},\mathfrak F(G_{\mathbb A})\bigr).
\]

\par}
\paragraph{Remark 2.1.4.} The space \(\mathfrak F(G,\GL(2,\Z))\) above is stable under product. On the other hand, \(D_{k-1}\otimes D_{\ell-1}\) contains the \(D_{k+\ell+2m-1}\), \(m\ge0\). For \(m=0\), this corresponds to the fact that the product \(fg\) of a holomorphic modular form of weight \(k\) and one of weight \(\ell\) is one of weight \(k+\ell\). For \(m=1\), in coordinates (1.5.2), one finds that
\[
\ell\,\frac{\partial f}{\partial z}\,g-kf\,\frac{\partial g}{\partial z}
\]
is a holomorphic modular form of weight \(k+\ell+2\), and so on. Likewise in the adelic framework.

\paragraph{Variant 2.1.5.} When, in the preceding considerations, \(D_{k-1}\) is replaced by another irreducible admissible representation endowed with a cyclic vector, one obtains the different types of nonholomorphic Mass modular forms.

\paragraph{(2.1.6)} The remark below is the analogue, for \(\GL(2,\R)\) and the \(D_{k-1}\), of the theorem of existence and uniqueness of the Kirillov model studied in the sections

\clearpage
\markboth{Authority 24 / printed 78}{}
\hypertarget{authority-24}{}
that follow.

\paragraph{Proposition 2.1.7.}{\itshape  Up to a scalar, there exists on \(G\) a unique function \(f_k\), holomorphic and of moderate growth, such that
\[
f_k(\lambda g)=\lambda^{-k}f_k(g),
\]
and such that
\[
f_k\!\left(g\begin{pmatrix}1&0\\u&1\end{pmatrix}\right)
=
e^{2\pi i u}f_k(g)
\qquad(u\in\R,\ g\in G).
\]
In the coordinate system (1.1.5.2), one has
\[
\tag{2.1.7.1}
f_k(\lambda,z)=
\begin{cases}
\lambda^{-k}e^{2\pi iz},& \text{for }\Im(z)>0,\\
0,& \text{for }\Im(z)<0.
\end{cases}
\]

\par}
The analogue of 2.1.7 for all infinite-dimensional irreducible admissible representations of both \(\GL(2,\R)\) and \(\GL(2,\C)\) appears in [9], §§5 and 6, or [15], Ch. VIII.

\subsection*{2.2. Kirillov model and new vector}

Let \(F\) be a nonarchimedean local field, let \(\psi\) be a nontrivial additive character of \(F\), and let
\[
\pi:\GL(2,F)\longrightarrow\GL(V)
\]
be an irreducible admissible representation of \(\GL(2,F)\). We shall also denote a uniformizer by \(\pi\); this should cause no confusion. The integer \(n(\psi)\) is the largest integer \(n\) such that \(\psi|_{\pi^{-n}\mathcal O}\) is trivial.

The statements that follow are purely algebraic. Thus one may take \(V\) to be a vector space over an algebraically closed field \(k\) of characteristic \(0\), with \(\psi\) taking values in the roots of unity of \(k\).

\paragraph{Proposition 2.2.1.}{\itshape  If \(V\) is finite-dimensional, then \(V\) is one-dimensional, and there exists a quasi-character \(\chi\) of \(F^*\) such that
\[
\pi(g)=\chi(\det g).
\]\par}

\clearpage
\markboth{Authority 25 / printed 79}{}
\hypertarget{authority-25}{}
Indeed, \(\Ker(\pi)\) is an open normal subgroup of \(\GL(2,F)\), hence contains \(\SL(2,F)\).

\paragraph{Theorem 2.2.2.}{\itshape  If \(V\) is infinite-dimensional, there exists a linear form
\[
L\ne0
\]
on \(V\), unique up to scalar, such that
\[
L\!\left(\pi\!\begin{pmatrix}1&0\\u&1\end{pmatrix}v\right)=\psi(u)L(v).
\]
\par}
A conceptual proof of this theorem, and a generalization to \(\GL(n)\), appear in [5], Th. 4 and 5.

Let \(W\) be the space of functions \(f\) on \(\GL(2,F)\) such that
\[
\tag{a}
f\!\left(\begin{pmatrix}1&0\\u&1\end{pmatrix}g\right)=\psi(u)f(g),
\]
\[
\tag{b}
\text{for }\varkappa\text{ in a sufficiently small neighborhood of }e,\quad f(g\varkappa)=f(g).
\]
The group \(\GL(2,F)\) acts on \(W\) by right translations.

By Frobenius reciprocity, Theorem 2.2.2 allows, in a unique way up to scalar, the realization of \(\pi\) in a subspace \(W(\pi)\) of \(W\): to \(v\in V\) one associates
\[
v(g)=L(\pi(g)v).
\]
\(W(\pi)\) is the \emph{Whittaker model} of \(\pi\). It depends only on the isomorphism class of \(\pi\). In the Whittaker model,
\[
L(v)=v(e).
\]

\paragraph{Proposition 2.2.3.}{\itshape  Let \(P\) be the subgroup
\[
\begin{pmatrix}*&0\\ *&*\end{pmatrix}
\]
of \(\GL(2,F)\). The restriction map from \(W(\pi)\) to the functions on \(P\) is injective.

\par}
Using the uniqueness of \(L\), one shows that the kernel consists of vectors invariant under
\[
\begin{pmatrix}1&0\\ *&1\end{pmatrix},
\]
and hence, by (2.2.5) below, under \(\SL(2,F)\).

See [5], Th. 6, for a generalization to \(\GL(n)\).

\clearpage
\markboth{Authority 26 / printed 80}{}
\hypertarget{authority-26}{}
% unit: deligne.d017.p026
Since \(\pi\) is irreducible, the center of \(\GL(2,F)\) acts by scalars: there exists a quasi-character \(\omega_\pi\) of \(\pi\) such that
\[
\pi\begin{pmatrix}a&0\\0&a\end{pmatrix}=\omega_\pi(a).
\]
In the preceding proposition, one may therefore replace \(P\) by the subgroup
\[
\begin{pmatrix}1&0\\0&*\end{pmatrix}
\]
of \(\GL(2,F)\). We denote by \(\calK(\pi)\) the set of functions on \(F^*\) of the form
\[
L\!\left(\begin{pmatrix}1&0\\0&a\end{pmatrix}v\right)\qquad (v\in V).
\]
The group \(\GL(2,F)\) acts on \(\calK(\pi)\) by transport of structure, from its action on \(V\). This is the \emph{Kirillov model} of \(\pi\). In this model, by construction,
\[
\tag{2.2.3.1}
\pi\begin{pmatrix}a&0\\b&d\end{pmatrix}f
=
\omega_\pi(a)\,\psi(a^{-1}bx)\,f(a^{-1}dx).
\]

\paragraph{Theorem 2.2.4.}{\itshape 
\begin{enumerate}[label=(\roman*)]
\item The functions in \(\calK(\pi)\) are locally constant, and their support is relatively compact in \(F\).
\item The space \(\mathcal S(F^*)\) of locally constant compactly supported functions in \(F^*\) is contained in \(\calK(\pi)\), and has finite codimension, \(0\), \(1\), or \(2\), in \(\calK(\pi)\).
\end{enumerate}
\par}
The difficult result is finite codimension. §2 of Godement--Jacquet [7] gives a conceptual proof of it.

\paragraph{Lemma 2.2.5.}{\itshape  Let \(G'_n\) be the subgroup of \(\SL(2,F)\) generated by the subgroups
\[
\begin{pmatrix}1&0\\ \OO&1\end{pmatrix}
\quad\text{and}\quad
\begin{pmatrix}1&\pi^n\OO\\0&1\end{pmatrix}.
\]
\begin{enumerate}[label=(\roman*)]
\item For \(n>0\), \(G'_n\) is the subgroup of \(\SL(2,\OO)\) formed by the
\[
\begin{pmatrix}a&c\\b&d\end{pmatrix}
\]
such that \(c\), \(a-1\), and \(d-1\) are \(0\) modulo \(\pi^n\).
\item For \(n=0\), \(G'_n=\SL(2,\OO)\).
\item For \(n<0\), \(G'_n=\SL(2,F)\).
\end{enumerate}
\par}
This is well known; we only indicate that (i) follows from the identity

\clearpage
\markboth{Authority 27 / printed 81}{}
\hypertarget{authority-27}{}
% unit: deligne.d017.p027
\[
\begin{pmatrix}1&-\dfrac{v}{1+uv}\\0&1\end{pmatrix}
\begin{pmatrix}1&0\\u&1\end{pmatrix}
\begin{pmatrix}1&v\\0&1\end{pmatrix}
\begin{pmatrix}1&0\\-\dfrac{u}{1+uv}&1\end{pmatrix}
=
\begin{pmatrix}\dfrac1{1+uv}&0\\0&1+uv\end{pmatrix}.
\]

\paragraph{Theorem 2.2.6.}{\itshape  Let \(\pi\) be an admissible irreducible infinite-dimensional representation of \(\GL(2,F)\). Let \(\alpha\) and \(\beta\) be two characters of \(\OO^*\) such that
\[
\alpha\beta=\omega_\pi.
\]
\begin{enumerate}[label=(\roman*)]
\item There exist nonzero vectors \(v\in V\) such that
\[
\begin{pmatrix}a&0\\b&d\end{pmatrix}v=\alpha(a)\beta(d)\,v
\qquad(a,d\in\OO^*,\ b\in\OO).
\]
\item Let \(n\) be the smallest integer such that there exists \(v\) as in (i), invariant under the
\[
\begin{pmatrix}1&c\\0&1\end{pmatrix}
\qquad(c\in\pi^n\OO).
\]
Then \(n\) is at least equal to the conductor of \(\alpha\beta^{-1}\), in particular \(n\ge0\).
\item Let \(G_k\) be the subgroup of \(\GL(2,F)\)
\[
G_k=
\begin{cases}
\GL(2,\OO),&k=0,\\[0.4em]
\left\{\begin{pmatrix}a&c\\b&d\end{pmatrix}\mid
a,d\in\OO^*,\ b\in\OO,\ c\in\pi^k\OO\right\},&k\ge1.
\end{cases}
\]
For \(k\) at least equal to the conductor of \(\alpha\beta^{-1}\), let \(\car(\alpha,\beta)\) be the character of \(G_k\) equal to \(\alpha(a)\beta(d)\) on
\[
\begin{pmatrix}a&c\\0&d\end{pmatrix}
\quad\text{and on}\quad
\begin{pmatrix}a&0\\b&d\end{pmatrix}.
\]
For \(k\ge n\), the space \(X_k\) of \(v\in V\) such that
\[
gv=\car(\alpha,\beta)(g)\,v\qquad(g\in G_k)
\]
has dimension \(k-n+1\), the direct sum of the lines
\[
\begin{pmatrix}1&0\\0&\pi^{-i}\end{pmatrix}X_n
\qquad(0\le i\le k-n).
\]
\end{enumerate}\par}

Take \(\psi\) such that \(n(\psi)=0\). In the Kirillov model, condition (i) then means that \(v\), identified with a function \(v(x)\) on \(F^*\), has support in \(\OO\) and satisfies
\[
v(dx)=\beta(d)v(x).
\]
Assertion (i) therefore follows from 2.2.3, \(\mathcal S(F^*)\subset\calK(\pi)\); assertion (ii) follows from 2.2.5.

The transformation
\[
\begin{pmatrix}1&0\\0&\pi^{i}\end{pmatrix}
\]
transforms vectors invariant under
\[
\begin{pmatrix}1&\pi^k\OO\\0&1\end{pmatrix}
\]
into vectors invariant under the group
\[
\begin{pmatrix}1&\pi^{k-i}\OO\\0&1\end{pmatrix},
\]
which is larger if \(i\ge0\). In the

\clearpage
\markboth{Authority 28 / printed 82}{}
\hypertarget{authority-28}{}
% unit: deligne.d017.p028
Kirillov model, it transforms functions supported in \(\pi^\ell\OO\) into functions supported in \(\pi^{\ell-i}\OO\), larger if \(i\ge0\). The definition of \(n\) gives
\[
X_n\cap
\begin{pmatrix}1&0\\0&\pi^i\end{pmatrix}X_{n+i-1}=0,
\]
that is,
\[
\begin{pmatrix}1&0\\0&\pi^{-i}\end{pmatrix}X_n\cap X_{n+i-1}=0.
\]
It remains to prove that
\[
\dim X_k\le \dim X_{k-1}+1.
\]
We shall prove that
\[
\dim\left(X_k\Big/
\begin{pmatrix}1&0\\0&\pi^{-1}\end{pmatrix}X_{k-1}\right)\le1.
\]
Indeed, let \(x,y\in X_k\). In the Kirillov model, the restrictions of \(x\) and \(y\) to \(\OO^*\) are necessarily proportional: there exists \((\lambda,\mu)\ne(0,0)\) such that \(\lambda x+\mu y\) has support in \(\pi\OO\). Hence
\[
\begin{pmatrix}1&0\\0&\pi\end{pmatrix}(\lambda x+\mu y)
\]
still satisfies condition (i) of the theorem, and so is in \(X_{k-1}\). This proves the assertion.

\paragraph{Definition 2.2.7.}{\itshape  Let \(\pi\) be admissible irreducible and infinite-dimensional, and apply 2.2.6 to the characters \(\omega_\pi\) and \(1\). The integer \(n\) of 2.2.6(ii) is the conductor of \(\pi\). The vector \(v\), unique up to a scalar, invariant under the
\[
\begin{pmatrix}a&c\\b&d\end{pmatrix}
\qquad
(d\in\OO^*,\ b\in\OO,\ c\equiv0(\pi^n),\ a\equiv1(\pi^n)),
\]
is the new vector of the representation.\par}

\paragraph{Remark 2.2.7.1.} In the Kirillov model, the new vector \(v\) is represented by a function \(\tilde v(x)\), supported in \(\OO\), depending only on the valuation of \(x\), and nonzero on \(\pi^{-n(\psi)}\OO^*\). This last point follows from the minimality of \(n\): consider
\[
\begin{pmatrix}1&0\\0&\pi\end{pmatrix}v.
\]
Standard arguments give the following corollaries.

\paragraph{Corollary 2.2.8.}{\itshape  For admissible irreducible representations of \(\GL(2,F)\), the field of definition of a representation \(\pi\), that is, the smallest field \(k_0\) such that \(\pi\)\par}

\clearpage
\markboth{Authority 29 / printed 83}{}
\hypertarget{authority-29}{}
{\itshape is isomorphic to \(\pi^\sigma\) for \(\sigma\in\Aut(\bar{k}/k_0)\),

is also the smallest field over which the representation can be written.\par}

It is the field of definition of the ideal in the group algebra which annihilates the new vector.

\paragraph{Corollary 2.2.9.}{\itshape  Let \(\alpha\) and \(\beta\) be two characters of \(\OO^*\), let \(n\) be the conductor of \(\alpha\beta^{-1}\), let \(G_n\) be as in 2.2.6, and let \(\calH\) be the algebra, for convolution product, of locally constant compactly supported functions on \(\GL(2,F)\) such that, for
\[
\gamma=\begin{pmatrix}a&c\\b&d\end{pmatrix}
\quad\text{and}\quad
\gamma'\in\begin{pmatrix}a'&c'\\b'&d'\end{pmatrix}
\]
in \(G_n\),
\[
f(\gamma g\gamma')=
\alpha(a)\beta(d)f(g)\alpha(a')^{-1}\beta(d')^{-1}.
\]
The algebra \(\calH\) is commutative.\par}

This corollary appears in Silberger, Proc. AMS 1959, p. 437--440. To deduce it from what precedes, one must use that the \(K\)-finite vectors of every unitary irreducible representation of \(\GL(2,F)\) form an admissible representation.

\subsection*{2.3. Kirillov model and Hecke operator}

The computations in this subsection are identical to the classical computations giving the action of Hecke operators on the expansions at the cusps of modular forms. The reason is explained in 2.5.

Let
\[
\pi:\GL(2,\Q_p)\longrightarrow\GL(V)
\]
be an admissible representation of \(\GL(2,\Q_p)\). Put
\[
K=\GL(2,\Z_p),\qquad
a=\begin{pmatrix}1&0\\0&p\end{pmatrix},
\]
and let \(\varphi\) be the characteristic function of the double class \(KaK\). This double class is the disjoint sum of the right classes
\[
\begin{pmatrix}1&0\\ i&p\end{pmatrix}K\quad(i\in\Z/p)
\qquad\text{and}\qquad
\begin{pmatrix}p&0\\0&1\end{pmatrix}K.
\]
The \emph{Hecke operator} \(T_p\) is convolution with \(\frac1p\varphi\). It is the composition of the projection
\[
\int\pi(k)v\,dk
\]
from \(V\) onto \(V^K\), and of an endomorphism of \(V^K\).

\clearpage
\markboth{Authority 30 / printed 84}{}
\hypertarget{authority-30}{}
% unit: deligne.d017.p030
Suppose that \(\pi\) is admissible irreducible of infinite dimension, and that \(V^K\ne0\). If \(v_0\) is the new vector of \(V\), then \(v_0\) generates \(V^K\), and
\[
T_pv_0=\lambda v_0
\]
for a scalar \(\lambda\). Let \(\psi\) be an additive character of \(\Q_p\), of conductor \(0\), i.e. \(\Z_p=\operatorname{Ker}(\psi)\). In the corresponding Kirillov model, \(v_0(x)\) is a function \(F_0\) of the valuation of \(x\), \(x\in F^*\), supported in \(\OO\). One has
\[
T_pv_0(x)=
\frac1p\sum_{i\in\Z/p}\psi(ix)v_0(px)
+\frac1p\,\omega_\pi(p)v_0(p^{-1}x)
\]
\[
=\ \text{restriction to }\OO\text{ of }
\left[v_0(px)+\frac1p\,\omega_\pi(p)v_0(p^{-1}x)\right].
\]
In terms of the function \(F_0\), assumed normalized so that
\[
F_0(0)=1,
\]
one obtains
\[
\lambda F_0(n)=F_0(n+1)+\frac1p\,\omega_\pi(p)F_0(n-1)
\qquad(n\ge0),
\]
that is,
\[
\sum_{n\ge0}F_0(n)t^n
=
\frac1{1-\lambda t+p^{-1}\omega_\pi(p)t^2}.
\]

\subsection*{2.4. Newforms of Atkin--Lehner [1]}

\paragraph{(2.4.1)} Let \(\mathcal A^0(k)\) be the space of holomorphic cuspidal modular forms of weight \(k\). It is an admissible representation of \(\GL(2,\Af)\), and it follows easily from the existence of the Petersson scalar product that it is an algebraic direct sum of admissible irreducible representations of \(\GL(2,\Af)\).

Let \(V\subset \mathcal A^0(k)\) be an irreducible subrepresentation. The center of \(\GL(2,\Af)\) acts on \(V\) by a quasi-character \(\omega\). The quasi-characters \(x^k\) of \(\R^*\) and \(\omega\) define a quasi-character of \(\A^*/\Q^*\). The quasi-character \(\omega\) is therefore determined by its restriction to \(\widehat{\Z}^{\,*}\), subject to satisfying
\[
\omega(-1)=(-1)^k.
\]
The representation \(\pi\) of \(\GL(2,\Af)\) on \(V\) is a restricted tensor product (0.2.3) of admissible irreducible representations of infinite dimension

\clearpage
\markboth{Authority 31 / printed 85}{}
\hypertarget{authority-31}{}
\(\pi_p\) of the local groups \(\GL(2,\Q_p)\):
\[
V=\bigotimes_p V_p.
\]
For each \(p\), let \(v_p\) be a new vector for \(\pi_p\), assuming that, for almost all \(p\), these \(v_p\) coincide with the vectors used to define the restricted tensor product (0.2.3). We say that
\[
v=\bigotimes v_p,
\]
well-defined up to a scalar factor, is the \emph{newform} of \(V\). A form \(f\in \mathcal A^0(k)\) is \emph{new} if it generates an irreducible subrepresentation of \(\mathcal A^0(k)\) and is a newform in it. Its \emph{conductor} is the product
\[
n=\prod_p p^{f_p},
\]
where \(f_p\) is the conductor of \(\pi_p\).

\paragraph{Scholium 2.4.2.}{\itshape  The map which associates to a newform \(f\) the representation of \(\GL(2,\Af)\) that it generates is a bijection from the set of newforms, up to scalar, to the set of irreducible subrepresentations of \(\mathcal A^0(k)\).\par}

\paragraph{(2.4.3)} Let \(f\) be a newform of conductor \(n\), and let \(\omega\) be the \emph{character} of \(f\), i.e. the character of \(\widehat{\Z}^{\,*}\) defined by the corresponding representation. By definition, for
\[
\gamma=\begin{pmatrix}a&c\\ b&d\end{pmatrix}\in\GL(2,\widehat{\Z}),
\qquad c\equiv0\pmod n,
\]
one has
\[
\tag{2.4.3.1}
\gamma f=\omega(a)f.
\]
By 1.2.5, a form \(f\) satisfying (2.4.3.1) is identified with a family
\[
f(z;\sigma)\qquad \bigl(\sigma\in\GL(2,\Z(n))\bigr)
\]
of functions on the Poincaré half-plane, satisfying (1.2.5.2), and
\[
f(z;\gamma^{-1}\sigma)=\omega(a)f(z;\sigma)
\]
for \(\gamma\) as above. This family is determined by the single function
\[
f(z)=f(z;1)
\]
on the Poincaré half-plane \(X\), satisfying

\clearpage
\markboth{Authority 32 / printed 86}{}
\hypertarget{authority-32}{}
\[
\tag{2.4.3.2}
f(z)=(cz+d)^{-k}\,\omega(a)^{-1}
f\!\left(\frac{az+b}{cz+d}\right)
\]
for
\[
\begin{pmatrix}a&c\\ b&d\end{pmatrix}\in\SL(2,\Z),
\qquad c\equiv0\pmod n
\]
(1.2.5).

If a form \(f\) satisfies (2.4.3.1), then, decomposing it according to the irreducible representations in \(\mathcal A^0(k)\), and applying (2.2.6), one sees that it is written uniquely as a linear combination of forms
\[
\begin{pmatrix}1&0\\0&\ell\end{pmatrix}f_{\ell,m},
\]
where \(\ell,m\) are integers \(\ge1\), with \(\ell m\mid n\), where \(\begin{pmatrix}1&0\\0&\ell\end{pmatrix}\) is considered as in \(\GL(2,\Af)\), and where \(f_{\ell,m}\) is a linear combination of newforms of conductor
\[
\frac{n}{\ell m}.
\]
In terms of functions on \(X\),
\[
\tag{2.4.3.3}
f(z)=\sum f_{\ell,m}(\ell z).
\]

\paragraph{Theorem 2.4.4.}{\itshape  Let \(k\) be an integer \(\ge1\), and let \(\omega\) be a character of \(\widehat{\Z}^{\,*}\) such that
\[
\omega(-1)=k.
\]
For \(f\) to be a linear combination of newforms of weight \(k\), character \(\omega\), and conductor equal to \(n\), it is necessary and sufficient that \(f\) satisfy (2.4.3.2), and that, for every newform \(g\) of weight \(k\), character \(\omega\), and conductor
\[
n/\ell m<n,
\]
one have
\[
\langle f(z),g(\ell z)\rangle=0
\]
for the Petersson scalar product.\par}

In the decomposition (2.4.3.3), one wants
\[
f_{\ell,m}=0\qquad(\ell m\ne1).
\]
Forms belonging to different representations are orthogonal for the Petersson scalar product, so that
\[
\langle f_{\ell,m}(\ell z),f_{\ell',m'}(\ell' z)\rangle=0
\qquad(\ell m\ne \ell' m').
\]
The assertion follows.

\paragraph{Variant 2.4.5.} Let \(\mathcal A^0\) be the space of functions \(f\) on
\[
G_{\A}/\GL(2,\Q)
\]
which are
\begin{enumerate}[label=(\alph*)]
\item \(C^\infty\), of moderate growth in \(g_\infty\);
\item invariant under an open subgroup of \(\GL(2,\Af)\);
\item cuspidal: \(f_P=0\) (1.3.2.1).
\end{enumerate}

By (2.1.2), Scholium 2.4.2 can also be stated by saying that the new holomorphic (cuspidal) modular forms are in bijection with the irreducible admissible sub-

\clearpage
\markboth{Authority 33 / printed 87}{}
\hypertarget{authority-33}{}
representations of \(\mathcal A^0\) of one of the types
\[
D_{k-1}\otimes\pi_f,
\]
where \(\pi_f\) is a representation of \(\GL(2,\Af)\) and \(k\ge1\). To obtain the others, one must also consider nonholomorphic forms.

\paragraph{(2.4.6)} In the rest of this paragraph, we explain how the representation of \(\GL(2,\Af)\) in \(\mathcal A^0(k)\) is written in classical terms.

Let \(\mathcal A_1^0(k)\) be the space of holomorphic functions on the Poincaré half-plane \(X\), such that
\[
f|_k\gamma=f
\]
for \(\gamma\) in a congruence subgroup of \(\SL(2,\Z)\), and which vanish at the cusps. The group \(\GL(2,\Q)^+\) of matrices with positive determinant acts on \(\mathcal A_1^0(k)\) by
\[
\begin{pmatrix}a&c\\ b&d\end{pmatrix}f
=
(cz+d)^{-k}f\!\left(\frac{az+b}{cz+d}\right).
\]
The stabilizer of each \(f\) contains a congruence subgroup of \(\SL(2,\Z)\), so that the action extends to an action of the completion of \(\GL(2,\Q)^+\) for the topology of the congruence subgroups of \(\SL(2,\Z)\). The completion is the subgroup
\[
\GL(2,\Af)_1=\GL(2,\Q)^+\cdot\SL(2,\widehat{\Z})
\]
of \(\GL(2,\Af)\), formed by the \(g\in\GL(2,\Af)\) of positive rational determinant.

\paragraph{Proposition 2.4.7.}{\itshape  The representation \(\mathcal A^0(k)\) of \(\GL(2,\Af)\) is induced by the representation \(\mathcal A_1^0(k)\) of \(\GL(2,\Af)_1\).\par}

The isomorphism with the induced representation is obtained by associating to \(f\in \mathcal A^0(k)\) the function
\[
f(z;\sigma)=f((z,1),\sigma)
\]
from \(\GL(2,\Af)\) into \(\mathcal A_1^0(k)\). Indeed,
\[
f(z;\sigma\gamma^{-1})=\gamma f(z;\sigma)
\]
for \(\gamma\in\GL(2,\Q)^+\), hence also in its completion.

\clearpage
\markboth{Authority 34 / printed 88}{}
\hypertarget{authority-34}{}
\subsection*{2.5. Expansion at the cusps}

\paragraph{(2.5.1)} Let \(\psi\) be the additive character of \(\A\), trivial on \(\Q\) and on \(\widehat{\Z}\), whose restriction to \(\R\) is
\[
e^{-2\pi iu}.
\]
For \(f\) a function on \(G_{\A}/\GL(2,\Q)\), put
\[
W_f(g)=
\int_{\A/\Q}
f\!\left(g\begin{pmatrix}1&0\\ u&1\end{pmatrix}\right)\psi(u)\,du,
\]
whence
\[
W_f\!\left(g\begin{pmatrix}1&0\\ u&1\end{pmatrix}\right)
=
\psi(u)^{-1}W_f(g).
\]
Since \(W_f\) is defined by an integration on the right, the map
\[
f\longmapsto W_f
\]
commutes with translations on the left. If \(f\) is a holomorphic modular form of weight \(k\), then \(W_f\) is still holomorphic in \(g_\infty\), \(C^\infty\), of moderate growth, and satisfies
\[
W_f(\lambda g)=\lambda^{-k}W_f(g)
\qquad(\lambda\in\C^*).
\]
Therefore \(W_f\) is the product of the function (2.1.7.1) on \(\GL(2,\R)\) by a function \(W'_f\) on \(\GL(2,\Af)\). We put
\[
K'_f(a)=
W'_f\!\left(\begin{pmatrix}1&0\\0&a^{-1}\end{pmatrix}\right).
\]
One has the defining formula
\[
\tag{2.5.1.1}
\int_{\A/\Q}
f\!\left(\bigl((z,1),
\begin{pmatrix}1&0\\0&a^{-1}\end{pmatrix}\bigr)
\begin{pmatrix}1&0\\u&1\end{pmatrix}\right)\psi(u)\,du
=
\exp(2\pi iz)\,K'_f(a).
\]

\paragraph{2.5.2.} Let \(V\subset \mathcal A^0(k)\) be an irreducible subrepresentation of \(\GL(2,\Af)\). Let
\[
V=\bigotimes V_p.
\]
From uniqueness of Whittaker models one deduces that
\[
f\longmapsto W_f(g^{-1})
\]
sends \(V\) onto the tensor product of the Whittaker models of the \(V_p\), namely the set of functions on \(\GL(2,\Af)\) which are linear combinations of functions
\[
\prod_p f_p(g_p),
\]
where \(f_p\) lies in the Whittaker model of \(V_p\), and where, for almost all \(p\), \(f_p\) is the function in this model which is \(\GL(2,\Z_p)\)-invariant and satisfies \(f_p(e)=1\). Hence
\[
f\longmapsto K'_f
\]
sends \(V\) onto the tensor product, in the same sense, of the Kirillov models of the \(V_p\).

\paragraph{(2.5.3)} Let \(f\) be a modular form of weight \(k\), of level \(N\), described as in (1.2.5) by
\[
f(z,\sigma)\qquad
(z\in X,\ \sigma\in\GL(2,\Z/N)).
\]
Put
\[
q=\exp(2\pi iz)
\]
and
\[
\tag{2.5.3.1}
f(z,\sigma)=\sum_n a(n,f,\sigma)q^n
\qquad(n\in\Q,\ n\ge0).
\]

\clearpage
\markboth{Authority 35 / printed 89}{}
\hypertarget{authority-35}{}
\paragraph{Proposition 2.5.3.2.}{\itshape  Let \(a_0\in\widehat{\Z}^{\,*}\) and \(n\in\Q\), \(n>0\). One has
\[
K'_f(a_0n)
=
n^{-k}\,
a\!\left(n,f,\begin{pmatrix}1&0\\0&a_0^{-1}\end{pmatrix}\bmod N\right).
\]\par}

Apply (2.5.1.1) and integrate \(u\) in
\[
[0,M]\times \frac1M\widehat{\Z}.
\]
For \(M\) sufficiently divisible, one obtains
\[
\exp(2\pi iz)K'_f(a_0n)
=
\frac1M\int_0^M
f\!\left((z+u,1),
\begin{pmatrix}1&0\\0&a_0^{-1}n^{-1}\end{pmatrix}\right)
\exp(-2\pi iu)\,du.
\]
Since
\[
f\!\left((z+u,1),
\begin{pmatrix}1&0\\0&a_0^{-1}n^{-1}\end{pmatrix}\right)
=
f\!\left((z+u,n),
\begin{pmatrix}1&0\\0&a_0^{-1}\end{pmatrix}\right)
\]
\[
=
n^{-k}
f\!\left(\left(\frac{z+u}{n},1\right),
\begin{pmatrix}1&0\\0&a_0^{-1}\end{pmatrix}\right),
\]
one still has
\[
\exp(2\pi iz)K'_f(a_0n)
=
\frac1M\int_0^M
n^{-k}f\!\left(\frac{z+u}{n},
\begin{pmatrix}1&0\\0&a_0^{-1}\end{pmatrix}\bmod N\right)
\exp(-2\pi iu)\,du,
\]
and 2.5.3.2 follows.

\paragraph{Corollary 2.5.4.}{\itshape  A holomorphic cuspidal modular form of weight \(k\), say \(f\), is entirely determined by the function
\[
K'_f\quad\text{on }\Af.
\]\par}

Suppose \(f\) has level \(N\), or a level dividing \(N\). Then \(f\) is determined by the
\[
f(z;\sigma)
\]
for \(\sigma\) of the form
\[
\begin{pmatrix}1&0\\0&a_0^{-1}\end{pmatrix},
\qquad
a_0\in(\Z/N\Z)^*,
\]
because these \(\sigma\) have arbitrary determinant. These are determined by their Fourier series expansion at the cusp \(i\infty\), i.e. by the
\[
a(n,f,\sigma)\qquad(n\in\Q).
\]
The \(a(n,f,\sigma)\), for \(n\le0\), are zero, because \(f\) has moderate growth and is cuspidal. For \(n>0\), they are therefore given by \(K'_f\).

Notice that, in this corollary, we have not assumed that \(f\) lies in a given irreducible subrepresentation of \(\mathcal A^0(k)\). This statement, being essentially global, is therefore not reduced to the local statement 2.2.3.

\paragraph{(2.5.5)} The preceding proof is translated as follows in adelic terms. First

\clearpage
\markboth{Authority 36 / printed 90}{}
\hypertarget{authority-36}{}
of all, by the Fourier inversion formula, for \(f\) cuspidal one has
\[
\tag{2.5.5.1}
f(g)=\sum_{n\in\Q^*}W_f\!\left(g\begin{pmatrix}1&0\\0&n\end{pmatrix}\right).
\]
Next observe that
\[
\GL(2,\R)\begin{pmatrix}1&0\\0&*\end{pmatrix}\GL(2,\Q)
\]
is dense in \(\GL(2,\A)\). Thus \(f\) is determined by the restriction of \(W_f\) to
\[
\GL(2,\R)\begin{pmatrix}1&0\\0&*\end{pmatrix},
\]
and, applying 2.1.7 as in 2.5.1, one obtains 2.5.4.

Let us recall the proof of the density result used. With the diagonal matrices available, and which may be put in front, one reduces to \(\SL_2\). For every place \(v\), \(\SL(2,\Q_v)\SL(2,\Q)\) is then dense in \(\SL(2,\A)\): for each \(v'\), \(\SL(2,\Q_{v'})\) is generated by its lower and upper unipotent subgroups, and this reduces to the analogous problem for the additive group.

\paragraph{Theorem 2.5.6.}{\itshape  Let \(S\) be a finite set of prime numbers. For \(p\notin S\), let \(\pi_p\) be an irreducible admissible representation of \(\GL(2,\Q_p)\). Then there exists at most one irreducible subrepresentation
\[
V\subset \mathcal A^0(k),\qquad V\simeq\bigotimes_p V_p,
\]
whose local components are isomorphic to the \(\pi_p\) for \(p\notin S\).\par}

Suppose such a representation exists, say \(V\simeq\bigotimes_pV_p\). For \(p\notin S\), let \(v_p\in V_p\) be the new vector of \(V_p\). For \(p\in S\), let \(v_p\in V_p\) be the vector which, in the Kirillov model of \(V_p\), is identified with the characteristic function of \(\OO^*\). The expression of \(v_p\) in the Kirillov model is determined by \(\pi_p\) for \(p\in S\); normalize \(v_p(x)\) so that \(v_p(1)=1\). It is given for \(p\in S\). In all cases, \(v_p(x)\) is a function of the valuation of \(x\), is supported in \(\OO\), and is equal to \(1\) at \(1\). For every integer \(n>0\), put
\[
a(n)=n^k\prod_p v_p(n).
\]
If \(f\in \mathcal A^0(k)\) is the form \(\bigotimes_pv_p\), then \(K'_f\) is known and, by 2.5.4, \(f\) is uniquely determined by the \((\pi_p)_{p\notin S}\). Since \(f\in V\), \(V\) is determined by these \(\pi_p\).

\clearpage
\markboth{Authority 37 / printed 91}{}
\hypertarget{authority-37}{}
One has
\[
f(z;\sigma)=\sum a(n)e^{2\pi inz},
\]
independent of \(\sigma\), for \(\sigma\) of the form
\[
\begin{pmatrix}1&0\\0&a\end{pmatrix},\qquad a\in\widehat{\Z}^{\,*}.
\]
In classical terms, the point of this proof is the possibility of modifying a modular form
\[
f=\sum a_nq^n
\]
into a modular form
\[
f=\sum a'_nq^n,
\]
with \(a'_n=0\) for \(p\mid n\) and \(p\in S\).

\paragraph{Proposition 2.5.7.}{\itshape  Let \(f\in \mathcal A^0(k)\) be a modular form of level \(N\), and
\[
f(z;\sigma)=\sum_{\lambda>0}a(\lambda,f,\sigma)q^\lambda.
\]
If \(f\) generates an irreducible subrepresentation \(V\subset \mathcal A^0(k)\) of \(\GL(2,\Af)\), and is a tensor product of vectors of the local components, then there exist a constant \(c\) and functions
\[
a_p(n,f,\sigma_p)\qquad (n\text{ an integer},\ \sigma_p\in\GL(2,\Z/p^{\nu_p(N)}))
\]
such that, for \(\lambda>0\),
\[
a(\lambda,f,\sigma)=c\prod_p a_p(\nu_p(\lambda),f,\sigma_p).
\]
Here \(\nu_p\) denotes the \(p\)-adic valuation and \(\sigma_p\) is the \(p\)-component of
\[
\sigma\in\GL(2,\Z/N)\simeq\prod_p\GL(2,\Z/N_p).
\]\par}

This says that \(W_f\) is a product. In particular, if
\[
f(z)=\sum a_ne^{2\pi inz}
\]
is a new form of weight \(k\), normalized so that \(a_1=1\), then:
\begin{enumerate}[label=\alph*)]
\item \(a_{nm}=a_na_m\) for \((n,m)=1\);
\item if \(p\) does not divide the conductor of \(f\), then \(f\) is an eigenvector of \(T_p\) (2.3). Put
\[
p^kT_pf=\lambda f,
\qquad \omega=\omega_k\varepsilon.
\]
One has (2.3) and (2.5.3.2)
\[
\sum_n a_{p^n}t^n=\frac{1}{1-\lambda t+p^{k-1}\varepsilon_p(p)t^2};
\]
\item \(f\) is uniquely determined by \(\omega\) and by the \(a_p\), for almost all \(p\); this is the translation of 2.5.6.
\end{enumerate}

\clearpage
\markboth{Authority 38 / printed 92}{}
\hypertarget{authority-38}{}
The analogue of b) for \(p\) dividing the conductor is treated in the same spirit in Casselman [3]; cf. 3.2.

\section*{3. Local class field theory and representations of \(\GL(2,F)\)}

\subsection*{3.1. Preliminaries}

\paragraph{3.1.1.} In this section we shall use the notation and results of §§2, 3, and 8 of [4] in this volume.

Let \(K\) be a local field. It is known what the \emph{Weil group} \(W(\overline K/K)\), relative to a separable closure \(\overline K\) of \(K\), is ([4], §2). For nonarchimedean \(K\), we introduced in [4], §8, a ``group'' \(W'(\overline K/K)\), which will enter only through the category of its representations. By definition, a \emph{representation} of \(W'(\overline K/K)\) on a finite-dimensional vector space \(V\) over a field \(k\) of characteristic \(0\) consists of:
\begin{enumerate}[label=(\alph*)]
\item a representation
\[
\sigma:W(\overline K/K)\longrightarrow\GL(V),
\]
trivial on an open subgroup of the inertia group \(I\);
\item a nilpotent endomorphism \(N\) of \(V\), such that
\[
wNw^{-1}=\omega_1(w)N\qquad(w\in W(\overline K/K)).
\]
\end{enumerate}
Recall that \(\omega_1\) is the unramified quasi-character giving the action of \(W(\overline K/K)\) on the roots of unity, and that \(\omega_n(x)=\omega(x)^n\). The geometric Frobenius elements \(F\) satisfy
\[
\omega_1(F)=\frac1q.
\]
Via the isomorphism of local class field theory, normalized so as to send geometric Frobenius elements to uniformizers, it is identified with the normalized absolute value.

Let \(E_\lambda\) be a finite extension of \(\Q_\ell\), \(\ell\ne p\). The set of isomorphism classes of representations of \(W'(\overline K/K)\) on \(E_\lambda\) is in canonical bijection with the set of isomorphism classes of \(\lambda\)-adic representations of \(W(\overline K/K)\). This is the reason for the existence of \(W'(\overline K/K)\).

\clearpage
\markboth{Authority 39 / printed 93}{}
\hypertarget{authority-39}{}
We shall have exclusively to consider \emph{\(F\)-semisimple representations} ([4], 8.6) of \(W'(\overline K/K)\), the Frobenius elements being semisimple. To unify the language, if \(K\) is archimedean and \(k=\C\), we define an \emph{\(F\)-semisimple representation} of \(W'(\overline K/K)\) to be a semisimple representation of the Lie group \(W(\overline K/K)\).

For the rest of this subsection, assume that \(K\) is nonarchimedean.

\paragraph{3.1.2.} Sums, quotients, and tensor products of representations of \(W'(\overline K/K)\) are defined in the evident way. Thus
\[
(\sigma',N')\otimes(\sigma'',N'')=(\sigma'\otimes\sigma'',N'\otimes1+1\otimes N'').
\]
Let \(\sprep(n)\) be the following representation \((\sigma,N)\) of \(W(\overline K/K)'\), on \(\Q\):
\begin{itemize}
\item the space of the representation is \(\Q^n\). We write \(e_i\), \(0\le i<n\), for the standard basis;
\item \(\sigma(w)e_i=\omega_i(w)e_i\);
\item \(Ne_i=e_{i+1}\) for \(0\le i<n\), and \(Ne_{n-1}=0\).
\end{itemize}

\paragraph{Proposition 3.1.3.}{\itshape 
\begin{enumerate}[label=(\roman*)]
\item The irreducible representations \((\sigma,N)\) of \(W'(\overline K/K)\) over \(k\) are the \((\sigma,0)\), with \(\sigma\) irreducible.
\item The indecomposable \(F\)-semisimple representations of \(W'(\overline K/K)\) over \(k\) are the \(\rho\otimes\sprep(n)\), with \(\rho\) irreducible.
\end{enumerate}\par}

Assertion (i) follows from the fact that \(\im(N)\) is a subrepresentation. Let \((\sigma,N)\) be an \(F\)-semisimple representation of \(W'(\overline K/K)\) on \(V\). The hypothesis implies that the representation \(\sigma\) is semisimple. For each isomorphism class \(\tau\) of representations of \(W(\overline K/K)\), let \(V(\tau)\) be the isotypic component of type \(\tau\) of \((V,\sigma)\), let \(R(\tau)\) be a representative of \(\tau\), let \(k(\tau)\) be the commutant of \(R(\tau)\), and let
\[
S(\tau)=\Hom_W(R(\tau),V).
\]
One has
\[
V=\bigoplus_\tau V(\tau)=\bigoplus R(\tau)\otimes_{k(\tau)}S(\tau).
\]

\clearpage
\markboth{Authority 40 / printed 94}{}
\hypertarget{authority-40}{}
The representation \((\sigma,N)\) is the direct sum of the subrepresentations
\[
\bigoplus_i V(\tau\otimes\omega_i),
\]
and is therefore equal to one of them. Let
\[
S=\bigoplus_i S(\tau\otimes\omega_i).
\]
Then
\[
V=R(\tau)\otimes_{k(\tau)}S.
\]
There is a unique way to make \(W'(\overline K/K)\) act on \(S\), by \((\sigma',N')\), so that this isomorphism is an isomorphism of representations. One then has
\[
\sigma'(w)|_{S(\tau\otimes\omega_i)}=\omega_i(w),
\qquad
N'S(\tau\otimes\omega_i)\subset S(\tau\otimes\omega_{i+1}).
\]
Choosing a basis of \(S\) adapted to the decomposition into the \(S(\tau\otimes\omega_i)\), and in which \(N'\) is in Jordan form, one obtains (ii).

\paragraph{Proposition 3.1.4.}{\itshape  Suppose \(k\) algebraically closed. Every irreducible representation \(\rho\) of \(W(\overline K/K)\), of dimension \(n<p\), is induced by a quasi-character of a subgroup of index \(n\) in \(W(\overline K/K)\).\par}

Fix \(p\), but not \(K\), and argue by induction, limited to \(p\), on \(n\). Recall that the dimension of every irreducible representation of a finite, or profinite, group divides the order of the group. Since \(n<p\), the restriction of the representation to \(P\subset I\) is a sum of one-dimensional representations. Let \(\chi\) be a character of \(P\) which occurs in \(\rho\), let \(V^\chi\) be the corresponding subspace of \(V\), and let \(W'\subset W(\overline K/K)\) be the stabilizer of \(\chi\), or of \(V^\chi\). There are two cases:
\begin{enumerate}[label=(\alph*)]
\item \(W'\ne W\). Then \(\rho\) is induced by the representation of \(W'\) on \(V^\chi\), to which the induction hypothesis is applied; \(W'\) is of the form \(W(\overline K/K)\).
\item \(W'=W\), so \(V^\chi=V\) and \(\rho(I)\) is commutative. There then exists a character of \(I\) which occurs in \(\rho\), and one argues as before: either one concludes by induction, or \(\rho(W(\overline K/K))\) is commutative, in which case \(n=1\).
\end{enumerate}

\paragraph{Example 3.1.5.} If \(p\ne2\), the complex \(F\)-semisimple representations of dimension two of \(W(\overline K/K)\) are the following:
\begin{enumerate}[label=(\alph*)]
\item the sum \([\lambda]\oplus[\mu]\) of two representations of dimension \(1\);
\item a representation \([\lambda]\otimes\sprep(2)\);
\item a representation induced by a quasi-character of a quadratic extension of \(K\).
\end{enumerate}

\clearpage
\markboth{Authority 41 / printed 95}{}
\hypertarget{authority-41}{}
\subsection*{3.2. Representations of \(\GL(2,K)\)}

This section is a fascicle of results. References for proofs will be given in 3.2.9.

\paragraph{3.2.1.} Let \(K\) be a local field. If \(\pi\) is a representation of \(\GL(2,K)\) and \(\omega\) is a quasi-character of \(K^*\), we shall write \(\pi\otimes\omega\) for the representation, realized in the same space, for which
\[
(\pi\otimes\omega)(g)=\pi(g)\,\omega(\det g),
\]
\(\check\pi\) for the admissible dual of \(\pi\), and \(\omega_\pi\) for the quasi-character of \(K^*\) such that
\[
\pi\!\left(\begin{smallmatrix}a&0\\0&a\end{smallmatrix}\right)=\omega_\pi(a).
\]
One has
\[
\tag{3.2.1.1}
\omega_{\pi\otimes\chi}=\omega_\pi\chi^2,
\qquad\text{and}
\]
\[
\tag{3.2.1.2}
\omega_{\check\pi}=\omega_\pi^{-1}.
\]

\paragraph{3.2.2.} Let \(\Phi\) be a nondegenerate bilinear form on \(K^2\), and let
\[
g\longmapsto g^\vee
\]
be the outer automorphism of \(\GL(2,K)\) such that
\[
\Phi(x,gy)=\Phi(g^\vee x,y).
\]
For \(\Phi\) alternating, one has
\[
g^\vee=\det(g)^{-1}g.
\]
The dual representation \(\check\pi\) is isomorphic to the representation \(g\mapsto\pi(g^\vee)\):
\[
\tag{3.2.2.1}
\check\pi\sim\bigl(g\longmapsto\pi(g^\vee)\bigr).
\]
In particular,
\[
\tag{3.2.2.2}
\check\pi\sim\pi\otimes\omega_\pi^{-1}.
\]
Assume a nontrivial additive character \(\psi\) has been chosen, and let \(\cK(\pi)\) be the space of the Kirillov model of \(\pi\). One has
\[
\tag{3.2.2.3}
\cK(\pi\otimes\omega)=\cK(\pi)\,\omega.
\]
Changing \(\psi\) does not modify the space \(\cK(\pi)\). In the Kirillov model, the duality between \(\pi\) and \(\pi\otimes\omega_\pi^{-1}\) is given by the following formula, for \(v_1\in\cK(\pi)\), \(v_2\in\cK(\pi\otimes\omega_\pi^{-1})\), and with \(v_1\) or \(v_2\) in \(\cS(F^*)\):

\clearpage
\markboth{Authority 42 / printed 96}{}
\hypertarget{authority-42}{}
\[
\tag{3.2.2.4}
\sprod{v_1}{v_2}=\int v_1(x)v_2(-x)\,d^*x.
\]

\paragraph{3.2.3.} When \(K\) is archimedean, or has residual characteristic \(\ne2\), and doubtless always, there exists a natural bijective correspondence between isomorphism classes of irreducible admissible representations of \(\GL(2,K)\) and isomorphism classes of \(F\)-semisimple representations of \(W'(\overline K/K)\). There are several natural ways of normalizing this isomorphism, changing \(\sigma\mapsto\pi(\sigma)\) into \(\sigma\mapsto\pi(\sigma)\otimes\omega_k\). We shall first assume that the base field of the representations is \(\C\), and describe the \emph{unitary correspondence}
\[
\sigma\longleftrightarrow\pi_u(\sigma),
\]
used implicitly in [7], [9], and [11]. It is characterized by the following properties, in fact already by (E).

\paragraph{(A) Functoriality.}
\[
\begin{array}{rlrl}
(1)&\pi_u(\sigma\otimes\omega)=\pi(\sigma)\otimes\omega,&&\text{(twist)},\\[0.35em]
(2)&\omega_{\pi_u(\sigma)}=\det\sigma,&&\text{(restriction to the center)},\\[0.35em]
(3)&\pi_u(\check\sigma)=\pi_u(\sigma)^\vee,&&\text{(automorphisms of }\GL(2)\text{, cf. 3.2.2.1).}
\end{array}
\]

\paragraph{(B) Discrete series.} In the nonarchimedean case,
\[
\begin{array}{rcl}
\pi_u(\sigma)\text{ supercuspidal}&\Longleftrightarrow&\sigma\text{ irreducible},\\
\pi_u(\sigma)\text{ in the discrete series}&\Longleftrightarrow&\sigma\text{ indecomposable}.
\end{array}
\]
In the archimedean case,
\[
\pi_u(\sigma)\text{ in the discrete series}\Longleftrightarrow
\sigma\text{ irreducible, i.e. indecomposable}.
\]
Recall that a representation is said to be in the discrete series, respectively supercuspidal, if one of its coefficients, hence all of them, is square-integrable, respectively compactly supported, modulo the center.

\clearpage
\markboth{Authority 43 / printed 97}{}
\hypertarget{authority-43}{}
\paragraph{(C) Induction.} Let \(\Ind(\chi_1,\chi_2)\) be the representation of \(\GL(2,K)\) induced by two quasi-characters \(\chi_1\) and \(\chi_2\) of \(K^*\), unitary induction. For \(K\) nonarchimedean, the constituents, i.e. the Jordan--Hölder quotients, of \(\Ind(\chi_1,\chi_2)\) are exactly the \(\pi_u(\sigma)\), for \(\sigma\) with semisimplification
\[
[\chi_1]\oplus[\chi_2].
\]
In all cases,
\[
\pi_u(\chi_1,\chi_2)={}_{\mathrm{dfn}}\,\pi_u([\chi_1]\oplus[\chi_2])
\]
is a constituent of \(\Ind(\chi_1,\chi_2)\).

\paragraph{(D) Degenerate representations.} The one-dimensional representations of \(\GL(2,F)\) are the \(\pi_u(\chi_1,\chi_2)\) for
\[
\chi_1\chi_2^{-1}=\omega_{\pm1}.
\]
For \(F\) real, and \(x\) the embedding of \(F\) in \(\C\), the irreducible representations of dimension \(n+1\) of \(\GL(2,F)\) are the \(\pi_u(\chi_1,\chi_2)\) for
\[
(\chi_1\chi_2^{-1})^{\pm1}=\omega_{-1}\,x^n.
\]
For \(F\) complex, and \(z\) and \(\overline z\) the embeddings of \(F\) in \(\C\), the finite-dimensional irreducible representations of \(\GL(2,F)\) are the \(\pi_u(\chi_1,\chi_2)\) for
\[
(\chi_1\chi_2^{-1})^{\pm1}=\omega_{-1}z^n\overline z^{-m}\qquad(n,m\ge0).
\]
In all cases, the trivial representation is
\[
\pi_u\bigl(\|x\|^{1/2},\|x\|^{-1/2}\bigr).
\]

\paragraph{(E) Tate functional equation.} Let \(\psi\) be a nontrivial additive character of \(K\), and \(dx\) the self-dual measure of \(K\) relative to \(\psi\). On the space \(\M_2(K)\) of matrices \(\left(\begin{smallmatrix}a&b\\ c&d\end{smallmatrix}\right)\), also denote by \(dx\) the measure \(da\,db\,dc\,dd\), and for \(f\) of Schwartz--Bruhat type on \(\M_2(K)\), put
\[
\widehat f(y)=\int f(x)\psi(\Tr(xy))\,dx.
\]
For \(\pi\) a representation of \(\GL(2,K)\) on \(V\), and \(f\in\cS(\M_2(K))\), denote by \(Z(f,\pi)\) the endomorphism of \(V\)
\[
Z(f,\pi)=\int f(g)\pi(g)\,d^*g,
\]
where \(d^*g\) is a Haar measure on \(\GL(2,F)\). This integral is defined by

\clearpage
\markboth{Authority 44 / printed 98}{}
\hypertarget{authority-44}{}
analytic continuation in the family of representations \(\pi\otimes\omega_s\), all realized in the same space.

The quotient
\[
\frac{Z(f,\pi_u(\sigma)\otimes\omega_s\otimes\omega_{1/2})}
     {L(\sigma\otimes\omega_s)}
\]
is a holomorphic function of \(s\), and a finite linear combination of coefficients of these endomorphisms, with \(f\) variable, is a zero-free function of \(s\). One has the functional equation
\[
\frac{{}^{t}Z\bigl(\widehat f,(\check{\pi}_u(\sigma)\otimes\omega_1)\otimes\omega_{1/2}\bigr)}
     {L(\check\sigma\otimes\omega_1)}
=
\varepsilon(\sigma,\psi,dx)
\frac{Z(f,\pi_u(\sigma)\otimes\omega_{1/2})}{L(\sigma)}.
\]

\paragraph{(F) Hecke functional equation.} Let \(\psi\) and \(dx\) be as in (E), and let \(\sigma\) be such that \(\pi_u(\sigma)\) is infinite-dimensional. For
\[
\varphi\in\cK(\pi(\sigma)\otimes\omega_{-1/2})=
\cK(\pi(\sigma))\omega_{-1/2}
\]
(3.2.2.3), the quotient
\[
\frac{\displaystyle\int_{K^*}\varphi(x)\omega_s(x)\,d^*x}
     {L(\sigma\otimes\omega_s)}
\]
is a holomorphic function of \(s\), the integral being defined by analytic continuation in \(s\). For suitable \(\varphi\), this function of \(s\) is zero-free. Put
\[
w=\begin{pmatrix}0&-1\\1&0\end{pmatrix}.
\]
For \(\varphi\in\cK(\pi(\sigma))\), one has the functional equation
\[
\frac{\displaystyle\int \omega_1(x)\omega_\pi^{-1}(x)\,w\varphi(x)\,\omega_{-1/2}(x)\,d^*x}
     {L(\omega_1\check\sigma)}
=
\varepsilon(\sigma,\psi,dx)
\frac{\displaystyle\int\varphi(x)\omega_{-1/2}(x)\,d^*x}{L(\sigma)},
\]

and hence, for every character \(\chi\) of \(F^*\),

\clearpage
\markboth{Authority 45 / printed 99}{}
\hypertarget{authority-45}{}
\[
\frac{\displaystyle\int \omega_1(x)\omega_\pi^{-1}(x)\chi^{-1}(x)\,w\varphi(x)\,\omega_{-1/2}(x)\,d^*x}
     {L(\omega_1\chi^{-1}\check\sigma)}
=
\varepsilon(\sigma\chi,\psi,dx)
\frac{\displaystyle\int\varphi(x)\chi(x)\omega_{-1/2}(x)\,d^*x}{L(\sigma\chi)}.
\]

\paragraph{(G) New vector.} Assume \(K\) nonarchimedean and \(\pi_u(\sigma)\) infinite-dimensional. The conductor of \(\pi_u(\sigma)\), defined in 2.2.7, then coincides with the conductor [4], 8.12.1 of \(\sigma\).

Let \(v_0\) be the new vector of \(\pi_u(\sigma)\). In the Kirillov model, \(v_0(x)\) is a function \(F'_0\) of the valuation of \(x\). Its support is contained in \(\pi^{-n(\psi)}\mathcal O\). Put
\[
F_0(n)=F'_0(n-n(\psi)),
\]
and normalize \(v_0\) so that \(F_0(0)=1\). One has
\[
\sum F_0(n)q^{-n/2}t^{nd}=L(\sigma\otimes\omega^t).
\]
Recall that \(q=p^d\) and that
\[
\omega^t(\pi)=t^d\qquad(\omega^{p^{-s}}=\omega_s).
\]

\paragraph{(3.2.4).} The chosen normalization of the correspondence \(\sigma\longleftrightarrow\pi_u(\sigma)\) has the advantage of giving rise to very simple formulas (A). From (A)(2) it also follows that \(\pi_u\) puts representations of \(\PGL(2,K)\) in bijective correspondence with \(F\)-semisimple representations of \(W'(\overline K/K)\) in \(\SL(2,\C)\).

In return, this normalization leads to scattering \(1/2\)'s through the formulas. In the nonarchimedean case, the only case where the problem arises, it is not invariant under automorphisms of \(\C\) if \(\sqrt q\) is not an integer.

\paragraph{(3.2.5).} Define the \emph{Tate correspondence} \(\sigma\longmapsto\pi_t(\sigma)\) by
\[
\pi_t(\sigma)=\pi_u(\sigma)\otimes\omega_{1/2}.
\]
Then
\[
\mathrm{A}(3)_t\qquad
\pi_t(\check\sigma\otimes\omega_1)=\check\pi_t(\sigma)\otimes\omega_2.
\]
Let \(\psi\) be a nontrivial additive character of \(K\), and \(dx\) the Haar measure

\clearpage
\markboth{Authority 46 / printed 100}{}
\hypertarget{authority-46}{}
self-dual on \(K\). We again write \(dx\) for the measure
\[
da\,db\,dc\,dd
\]
on \(M_2(K)\), and define \(\widehat f\) by
\[
\widehat f(x)=\int f(y)\,\psi(\Tr(xy))\,dy.
\]
The Tate functional equation takes the rationalized form
\[
\mathrm{(E)}_t\qquad
\frac{{}^{t}Z(\widehat f,\check\pi_t(\sigma)\otimes\omega_2)}
     {L(\check\sigma\otimes\omega_1)}
=
\varepsilon(\sigma,\psi,dx)\,
\frac{Z(f,\pi_t(\sigma))}{L(\sigma)}.
\]
For arbitrary \(dx\), multiply the right-hand side by \(dx/dx'\).

\paragraph{(C) Induction.} This becomes: \(\pi_t(\chi_1,\chi_2)\) occurs in the representation of \(\GL(2,K)\), by right translation, on the space of functions on \(\GL(2,F)\), right \(K\)-finite, such that
\[
\tag{5.2.3.1}
f\!\left(\begin{pmatrix}a&b\\0&d\end{pmatrix}g\right)
=\chi_1\omega_1(a)\,\chi_2(d)\,f(g).
\]

\paragraph{(D) Degenerate representations.} This becomes:
\[
\text{a) the unit representation is }\pi_t(\omega_{-1},1),
\]
\[
\text{b) for }K=\R\text{ or }\C,\text{ and }V\text{ the evident two-dimensional representation,}
\]
\[
\Sym^n(V)=\pi_t(\omega_{-1}x^n,1)\qquad\text{(real case)},
\]
\[
\Sym^n(V)\otimes\Sym^m(\overline V)
=\pi_t(\omega_{-1}z^n\overline z^{-m},1)
\qquad\text{(complex case)}.
\]

\paragraph{(5.2.6).} Define the \emph{Hecke correspondence} by
\[
\pi_h(\sigma)=\pi_u(\sigma)\otimes\omega_{-1/2}
=\pi_t(\sigma)\otimes\omega_{-1}.
\]
This time, \(\pi_h(\chi_1,\chi_2)\) occurs in the space of functions satisfying
\[
f\!\left(\begin{pmatrix}a&b\\0&d\end{pmatrix}\right)
=\chi_1(a)\chi_2(d)\omega_{-1}(d)f(g),
\]
and, with the notation of (D) above, one has

\clearpage
\markboth{Authority 47 / printed 101}{}
\hypertarget{authority-47}{}
\[
\mathrm{(D)}_h\qquad
1=\pi_h(1,\omega_1),
\]
\[
\Sym^n(V)=\pi_h(x^n,\omega_1)\qquad(K=\R),
\]
\[
\Sym^n(V)\otimes\Sym^m(\overline V)
=\pi_h(z^n\overline z^{-m},\omega_1)\qquad(K=\C).
\]

The Hecke functional equation takes the following simple form. Let \(\psi\) be a nontrivial additive character of \(K\), and let \(dx\) be self-dual; for arbitrary \(dx\), multiply the right-hand side by \(dx'/dx\). For \(\varphi\) in the Kirillov model of \(\pi_h(\sigma)\),
\[
\mathrm{(F)}_h\qquad
\frac{\displaystyle\int \omega_1(x)\omega_\pi^{-1}(x)\chi^{-1}(x)\,w\varphi(x)\,d^*x}
     {L(\omega_1\chi^{-1}\check\sigma)}
=
\varepsilon(\sigma\chi,\psi,dx)
\frac{\displaystyle\int \chi(x)\varphi(x)\,d^*x}{L(\chi\sigma)}.
\]
With the notation of (G), the new vector
\[
v_0(x)=F_0(v(x)+n(\psi))
\]
of \(\calK(\pi_h(\sigma))\) satisfies
\[
\mathrm{(G)}_h\qquad
\sum F_0(n)t^{nd}=L(\sigma\otimes\omega^t).
\]

\paragraph{3.2.7.} For \(K\) nonarchimedean, the correspondence
\[
\sigma\longleftrightarrow\pi_t(\sigma)\qquad(\text{respectively }\sigma\longleftrightarrow\pi_h(\sigma))
\]
is invariant under automorphisms of \(\C\). Hence, by (2.2.8), if the isomorphism class of \(\sigma\) is defined over a field \(K\) of characteristic \(0\), then \(\pi_t(\sigma)\), respectively \(\pi_h(\sigma)\), is realized over \(K\). Thus, to every \(\ell\)-adic representation
\[
\sigma:W(\overline K/K)\longrightarrow\GL(2,\Q_\ell)
\]
there corresponds an isomorphism class of representations \(\pi_t(\sigma)\), respectively \(\pi_h(\sigma)\), of \(\GL(2,K)\) over \(\Q_\ell\).

In the conjectural correspondence between isogeny classes of elliptic curves over \(\Q\) and new forms of weight \(2\), the conjecture is that the \(\ell\)-adic representation
\[
H_1(E,\Z_\ell)=T_\ell(E),
\]
restricted to the decomposition group at \(p\), corresponds under the Hecke correspondence to the factor at \(p\) of the representation of

\clearpage
\markboth{Authority 48 / printed 102}{}
\hypertarget{authority-48}{}
\(\GL(2,\A)\) associated with the newform.

\paragraph{3.2.8.} Suppose \(K\not\simeq\C\), and let \(H\) be ``the'' quaternion division algebra over \(K\). Since every automorphism of \(H\) is inner, the group \(H^*\) depends, up to inner automorphism, only on \(K\). The groups \(\GL(2,F)\) and \(H^*\) are inner forms of one another. Hence:

\begin{enumerate}[label=(\alph*)]
\item One can compare conjugacy classes in \(\GL(2,F)\) and \(H^*\). One finds that the set of conjugacy classes in \(H^*\) is identified with the set of elliptic or central conjugacy classes in \(\GL(2,F)\).
\item \(\bigwedge^4\Lie(\GL(2,F))\) and \(\bigwedge^4\Lie(H^*)\) are canonically isomorphic. Passing to absolute values, one obtains a correspondence between Haar measures on \(\GL(2,F)\) and Haar measures on \(H^*\).
\end{enumerate}
The notation (3.2.1) is transferred to \(H^*\) in the evident way. There is a natural bijective correspondence
\[
\pi\longmapsto \pi'
\]
between isomorphism classes of irreducible admissible representations, of finite dimension, of \(H^*\), and isomorphism classes of irreducible admissible representations in the discrete series of \(\GL(2,F)\). This correspondence is characterized by the following properties.

\paragraph{(H).} The analogue of (3.2.2.1), (3.2.2.2) is true for representations of \(H^*\). One has
\[
\mathrm{(H1)}\qquad (\pi\otimes\omega)'=\pi'\otimes\omega,
\]
\[
\mathrm{(H2)}\qquad \omega_\pi=\omega_{\pi'},
\]
\[
\mathrm{(H3)}\qquad (\pi^\vee)'=(\pi')^\vee.
\]

\paragraph{(I).} The conductor of \(\pi'\) is the smallest integer \(n\) such that \(\pi(x)\) is trivial for \(x-1\) of valuation \(\ge n\) in \(H\).

\paragraph{(J).} On elliptic elements, cf. (a), the character of \(\pi\) is the negative of the character of \(\pi'\). The formal degrees of \(\pi\) and \(\pi'\) are equal, cf. (b).

\paragraph{(K).} Up to a change of sign, \(\pi\) and \(\pi'\) satisfy the same Tate functional equation: for Euler factors \(L(\pi,s)\) and suitable ``constants'',

\clearpage
\markboth{Authority 49 / printed 103}{}
\hypertarget{authority-49}{}
one has
\[
\frac{{}^{t}Z(\widehat f,\pi'^{\vee}\otimes\omega_{2-s})}{L(\pi^{\vee},2-s)}
=
\varepsilon(s,\pi,\psi)\,
\frac{Z(f,\pi'\otimes\omega_s)}{L(\pi,s)},
\]
and
\[
\frac{{}^{t}Z(\widehat g,\pi^{\vee}\otimes\omega_{2-s})}{L(\pi^{\vee},2-s)}
=-\varepsilon(s,\pi,\psi)\,
\frac{Z(g,\pi\otimes\omega_s)}{L(\pi,s)},
\]
with quotients \(Z/L\) holomorphic in \(s\), and with some linear combination sometimes nonvanishing, as in (E). Here \(\psi\) denotes an additive character of \(K\), and the Fourier transforms are taken relative to a Haar measure self-dual for \(\psi\Tr(xy)\), or \(\psi\Tr\red(xy)\).

\paragraph{3.2.9.} Here are references for the proofs. Unless expressly stated otherwise, residual characteristic \(2\) is not excluded.

\paragraph{3.2.9.1.} A conceptual proof of (3.2.2.1) appears in [5], Th. 2. The proof in loc. cit. is valid for \(\GL(n)\). Formula (3.2.2.4) is the only one compatible with the invariance of \(\langle\quad\rangle\) under the lower triangular group. A direct proof is also given in [6], which also treats the archimedean case.

\paragraph{3.2.9.2.} The existence of a Tate functional equation is proved a priori in [7] for irreducible admissible representations of the multiplicative group of any simple algebra.

\paragraph{3.2.9.3.} The equivalence between (E) and (F), the Tate and Hecke functional equations, is proved in [9], §13.

\paragraph{3.2.9.4.} The results (3.2.8), (H), (J), (K) make up §§1 and 15 of [9]. (I) was pointed out to me by Weil, and follows rather easily from (G) and (K).

\clearpage
\markboth{Authority 50 / printed 104}{}
\hypertarget{authority-50}{}
\paragraph{3.2.9.5.} The proof of (G) is given in [3].

\paragraph{3.2.9.6.} The dictionary \(\sigma\longleftrightarrow\pi(\sigma)\) for representations \(\pi\) which are constituents of an induced representation, all of them for \(K\) archimedean, and the proof of (A), (B), (C), (D), (F) for these are given, for \(K=\R\), in [9], §5, for \(K=\C\) in [9], §6, and for \(K\) nonarchimedean in [9], §3. See also [6].

\paragraph{3.2.9.7.} The dictionary \(\sigma\longrightarrow\pi(\sigma)\) for \(\sigma\) induced by a character of a quadratic extension, and the proof of (F), are given in [9], 4.7. I have not checked the injectivity of \(\sigma\longrightarrow\pi(\sigma)\). I was told that it could be read from the formulas giving the character. Jacquet knows how to deduce it from his theory [9 bis].

\paragraph{3.2.9.8.} By 3.1.5, if \(p\ne2\), all representations \(\sigma\) have now been considered. By [7], §2, all non-supercuspidal representations \(\pi\) have been considered as well. The forgotten representations \(\pi\) therefore have nonzero mass in the Plancherel formula. If \(p\ne2\), completeness of the dictionary is therefore read from the Plancherel formula in its explicit form; see for instance [2].

\paragraph{3.2.9.8.} Langlands verified, using (3.2.8) and (I), that, for \(\Q_2\), there exist further representations \(\pi\) than those considered up to now. There also exist further \(\sigma\)'s.

In equal characteristic, even characteristic \(2\), [9], §12 proves by a global method that to every \(F\)-semisimple representation \(\sigma\) of \(W'(\overline K/K)\) there corresponds a unique representation \(\pi(\sigma)\) satisfying (F), and (E). The proof uses Artin's conjecture, known for function fields.

\clearpage
\markboth{Authority 51 / printed 105}{}
\hypertarget{authority-51}{}
\begin{center}
\textbf{Bibliography}
\end{center}

\noindent\textbf{[1]}\quad A. O. L. Atkin and J. Lehner -- Hecke operators on \(\Gamma_0(m)\). Math. Ann. \underline{185} (1970) 134--160.\par\medskip

\noindent\textbf{[2]}\quad P. Cartier\par\medskip

\noindent\textbf{[3]}\quad W. Casselman -- On some results of Atkin and Lehner. Math. Ann. \underline{201} (1973), p. 301--314.\par\medskip

\noindent\textbf{[4]}\quad P. Deligne -- Les constantes des équations fonctionnelles des fonctions \(L\). This volume.\par\medskip

\noindent\textbf{[5]}\quad I. M. Gel'fand and D. A. Kajdan -- Representations of the group \(\GL(n,K)\) where \(K\) is a local field in: Lie groups and their representations (Budapest, 1971), p. 91--118 (Halsted).\par\medskip

\noindent\textbf{[6]}\quad R. Godement -- Notes on Jacquet--Langlands theory. IAS Princeton 1970.\par\medskip

\noindent\textbf{[7]}\quad R. Godement and H. Jacquet -- Zêta functions of simple algebras -- Lecture Notes in Math. \underline{260} -- Springer Verlag 1972.\par\medskip

\noindent\textbf{[8]}\quad E. Hecke -- Mathematische Werke -- Vandenhoeck \& Reprecht. Göttingen 1970.\par\medskip

\noindent\textbf{[9]}\quad H. Jacquet and R. P. Langlands -- Automorphic forms on \(\GL(2)\). Lecture Notes in Math. \underline{114}, Springer Verlag.\par\medskip

\noindent\textbf{[9 bis]}\quad H. Jacquet. Automorphic forms on \(\GL(2)\). II. Lecture Notes \underline{278}.\par\medskip

\noindent\textbf{[10]}\quad J. Labesse -- \(L\)-undistinguishable representations and trace formula for \(\SL(2)\) in: Lie groups and their representations (Budapest, 1971), p. 331--338\par\medskip

\noindent\textbf{[11]}\quad R. P. Langlands -- Euler Products. James K. Whittemore Lecture in Math. (Halsted). Yale University, 1967.\par\medskip

\noindent\textbf{[12]}\quad R. P. Langlands -- Problems in the theory of Automorphic Forms -- in: Conference on Modern analysis and applications, Washington 1969, ), Lecture Notes in Math. 170, Springer Verlag.\par\medskip

\noindent\textbf{[13]}\quad A. Robert -- Formes automorphes sur \(\GL(2)\) (Travaux de H. Jacquet et R. P. Langlands). Séminaire Bourbaki 415, juin 1972. Lecture Notes in Math. \underline{317} Springer Verlag.\par\medskip

\noindent\textbf{[14]}\quad A. J. Silberger -- \(\mathrm{PGL}(2)\) over the \(p\)-adics, its representations, spherical functions and Fourier analysis. Lecture Notes in Math. 166, Springer Verlag 1970.\par\medskip

\noindent\textbf{[15]}\quad A. Weil -- Dirichlet series and automorphic forms. Lecture Notes in Math. 181, Springer Verlag 1971.\par\medskip
\end{document}
